\documentclass[10pt,a4paper]{article}
\usepackage[utf8]{inputenc}
\usepackage[T1]{fontenc}
\usepackage{lmodern}
\usepackage[french]{babel}
\usepackage[strings]{underscore}
\usepackage[a4paper,landscape,margin=1.4cm,headheight=15pt]{geometry}
\usepackage{array}
\usepackage{booktabs}
\usepackage{longtable}
\usepackage{xcolor}
\usepackage{colortbl}
\usepackage{enumitem}
\usepackage{fancyhdr}
\usepackage{titlesec}
\usepackage[hidelinks,bookmarks=true]{hyperref}
\definecolor{cBlock}{HTML}{B00020}
\definecolor{cHigh}{HTML}{C75300}
\definecolor{cMed}{HTML}{8A6D00}
\definecolor{cLow}{HTML}{5A5A5A}
\definecolor{cOk}{HTML}{2E7D32}
\definecolor{cAccent}{HTML}{1F4E79}
\definecolor{rowalt}{HTML}{EEF2F6}
\definecolor{boxbg}{HTML}{F4F6F8}
\definecolor{rule}{HTML}{B9C2CC}
\pagestyle{fancy}\fancyhf{}
\renewcommand{\headrulewidth}{0.4pt}\renewcommand{\footrulewidth}{0.2pt}
\fancyhead[L]{\small\textcolor{cAccent}{\textbf{Standards de stockage par OS \& version}}}
\fancyhead[R]{\small\textcolor{cLow}{Recherche sourc\'ee}}
\fancyfoot[L]{\small\textcolor{cLow}{À v\'erifier contre les sources cit\'ees --- les standards \'evoluent}}
\fancyfoot[R]{\small\thepage}
\titleformat{\section}{\color{cAccent}\Large\bfseries}{\thesection}{0.6em}{}
\titleformat{\subsection}{\color{cAccent}\large\bfseries}{\thesubsection}{0.6em}{}
\setlist{nosep,leftmargin=1.4em}
\renewcommand{\arraystretch}{1.2}
\setlength{\tabcolsep}{3pt}
\newcommand{\callout}[2]{\begin{center}\fcolorbox{rule}{boxbg}{\begin{minipage}{0.97\linewidth}\textbf{\textcolor{cAccent}{#1}}\\[2pt]#2\end{minipage}}\end{center}}
\newcolumntype{M}{>{\scriptsize\ttfamily\raggedright\arraybackslash}p{3.0cm}}
\newcolumntype{R}{>{\scriptsize\raggedright\arraybackslash}p{1.7cm}}
\newcolumntype{S}{>{\scriptsize\centering\arraybackslash}p{1.1cm}}
\newcolumntype{Z}{>{\scriptsize\raggedright\arraybackslash}p{1.9cm}}
\newcolumntype{T}{>{\scriptsize\raggedright\arraybackslash}p{1.4cm}}
\newcolumntype{O}{>{\scriptsize\ttfamily\raggedright\arraybackslash}p{4.6cm}}
\newcolumntype{C}{>{\scriptsize\centering\arraybackslash}p{1.3cm}}
\newcolumntype{I}{>{\scriptsize\centering\arraybackslash}p{0.9cm}}
\newcolumntype{N}{>{\scriptsize\raggedright\arraybackslash}p{4.4cm}}
\begin{document}
\begin{titlepage}\centering\vspace*{2cm}
{\color{cAccent}\rule{\linewidth}{1.4pt}}\\[0.5cm]
{\Huge\bfseries\color{cAccent} Standards de stockage par OS \& version}\\[0.3cm]
{\LARGE Partitionnement, systèmes de fichiers, options de montage}\\[0.2cm]
{\large Compilation sourc\'ee --- AIX, Linux (RHEL/Debian/Ubuntu/SUSE/Oracle), Windows Server, FHS/CIS}\\[0.4cm]
{\color{cAccent}\rule{\linewidth}{1.4pt}}\\[1cm]
\begin{minipage}{0.8\linewidth}\centering\itshape Base de r\'ef\'erence normalisable : pour chaque OS et version, le layout recommand\'e/par d\'efaut (points de montage / volumes), s\'eparation, tailles, type de FS, options de montage (CIS) et classification --- avec les sources officielles cit\'ees.\end{minipage}\\[1.4cm]
{\large Date : 12 juin 2026}\vfill
{\color{cLow}\small Document de travail --- usage interne}\end{titlepage}
\tableofcontents\thispagestyle{fancy}\clearpage
\section*{Avertissement \& m\'ethode}\addcontentsline{toc}{section}{Avertissement \& m\'ethode}
\callout{À lire avant usage}{Les chiffres ci-dessous proviennent de la documentation officielle des \'editeurs et des CIS Benchmarks (sources cit\'ees en fin de document). \textbf{Les standards \'evoluent par version} et certaines tailles sont des minima d'installation ou des recommandations, pas des obligations strictes. \textbf{Toujours v\'erifier contre la source cit\'ee} avant de figer un \og golden image\fg{}. Les colonnes sont normalis\'ees pour alimenter directement le modèle de r\'ef\'erence du projet (\texttt{reference/}).}
\medskip\noindent\textbf{L\'egende des colonnes :} \emph{Sep.} = filesystem s\'epar\'e exig\'e ; \emph{Req.} = \textcolor{cBlock}{mandatory}/\textcolor{cHigh}{recommended}/\textcolor{cMed}{conditional}/\textcolor{cLow}{optional} ; \emph{Opt.} = options de montage (surtout CIS) ; \emph{Cls} = system/data/boot/swap ; \emph{Src} = renvoi à la source (section Sources). Tailles dans l'unit\'e de la source.
\clearpage
\section{RHEL family}
\noindent\small\textit{This reference model covers official Red Hat Enterprise Linux recommended storage partitioning standards for RHEL 7, 8, and 9 based on: (1) Red Hat official Installation Guides for each version, (2) Red Hat Security Hardening and Security Guide documentation, (3) CIS Benchmarks for RHEL 7, 8, and 9 which specify secure mount options (nodev, nosuid, noexec), and (4) Filesystem Hierarchy Standard (FHS 3.0) principles. RHEL 7.3 increased /boot default from 500 MB to 1 GB. All versions use XFS as the default filesystem for / and data partitions, but /boot is restricted to ext4 (not XFS). Mount options recommended by CIS Benchmarks (nodev, nosuid, noexec) are included where security controls are appropriate. Sizes follow official Red Hat documentation (250 MiB minimum for /boot in earlier RHEL 7, 1 GiB default from 7.3 onward; 5 GiB minimum for /var in RHEL 8/9; 3 GiB for RHEL 7). Swap sizing guidance reflects modern Red Hat recommendations that swap should be sized by workload, not RAM amount. Note: RHEL 7 sources marked retrieved=false indicate documentation exists but content was behind authentication walls; searches confirmed information present in official documentation. RHEL 8/9 sources similarly authenticated but information confirmed via web search results citing official sources. Rocky Linux, Oracle Linux, and AlmaLinux are binary-compatible RHEL rebuilds and follow identical partition recommendations as their RHEL base versions. This data is intended for a normalized reference model and LaTeX PDF documentation."}\normalsize\par\medskip
\subsection{RHEL 7}
\noindent\small \textbf{FS par d\'efaut :} \texttt{XFS} \quad\textbullet\quad \textbf{Swap/paging :} Automatic partitioning limits swap partition to maximum 10\% of total disk or 128 GB, whichever is smaller. Modern systems with hundreds of GB of RAM should consider swap as a function of system memory workload, not system memory. Minimum 256 MB for systems with 1 GB or less RAM. For hibernation, swap space must be at least equal to RAM size.\normalsize\par\smallskip
\noindent\small\textit{RHEL 7.3 increased /boot default from 500 MB to 1 GB. /boot must always be on a physical partition, not LVM. Default filesystem for / and /boot is ext4 (XFS cannot be used for /boot). Default partitioning scheme uses LVM.}\normalsize\par\smallskip
{\renewcommand{\arraystretch}{1.15}\begin{longtable}{M R S Z Z T O C I N}\toprule\textbf{\normalfont\scriptsize Point de montage / volume} & \textbf{\normalfont\scriptsize Req.} & \textbf{\normalfont\scriptsize Sep.} & \textbf{\normalfont\scriptsize Min} & \textbf{\normalfont\scriptsize Max} & \textbf{\normalfont\scriptsize FS} & \textbf{\normalfont\scriptsize Options de montage} & \textbf{\normalfont\scriptsize Cls} & \textbf{\normalfont\scriptsize Src} & \textbf{\normalfont\scriptsize Notes}\\\midrule\endhead
/ & \textcolor{cBlock}{mandatory} & non & 3.0 GiB & unbounded & XFS & defaults,relatime & system & rhel7-install-guide & Root filesystem. 3.0 GiB allows minimal installation; 5.0 GiB allows full installation with packages.\\
\rowcolor{rowalt}
/boot & \textcolor{cBlock}{mandatory} & \textcolor{cOk}{oui} & 250 MiB & unbounded & ext4 & defaults,relatime & boot & rhel7-install-guide & Must be on physical partition, not LVM. Default size is 1 GB (increased from 500 MB in RHEL 7.3). Each kernel requires approximately 30-60 MiB.\\
/boot/efi & \textcolor{cMed}{conditional} & \textcolor{cOk}{oui} & 50 MiB & unbounded & vfat & defaults,umask=0077,shortname=winnt & boot & rhel7-install-guide & Required for UEFI systems on AMD64, Intel 64, and 64-bit ARM. Recommended minimum size is 200 MiB. Must not be encrypted.\\
\rowcolor{rowalt}
/home & \textcolor{cHigh}{recommended} & \textcolor{cOk}{oui} & 1 GiB & unbounded & XFS & defaults,nodev,nosuid,relatime & data & rhel7-security-guide,cis-rhel7 & Separate /home partition protects user data during system upgrades. Prevents / partition from filling and destabilizing system.\\
/var & \textcolor{cHigh}{recommended} & \textcolor{cOk}{oui} & 3.0 GiB & unbounded & XFS & defaults,nodev,nosuid,relatime & data & rhel7-install-guide,rhel7-security-guide & Holds logs, spool, temporary files, and application data. Must be large enough for package updates. Prevents / from filling due to log growth.\\
\rowcolor{rowalt}
/var/log & \textcolor{cHigh}{recommended} & conditional & 2.0 GiB & unbounded & XFS & defaults,nodev,noexec,nosuid,relatime & data & cis-rhel7,rhel7-security-guide & Separate partition prevents log growth from affecting system stability. CIS Benchmark recommends nodev, noexec, nosuid mount options.\\
/var/log/audit & \textcolor{cHigh}{recommended} & conditional & 1.0 GiB & unbounded & XFS & defaults,nodev,noexec,nosuid,relatime & data & cis-rhel7,rhel7-security-guide & Should reside on separate mount point to prevent other processes from consuming audit log space and provide accurate remaining space detection. CIS Benchmark recommends nodev, noexec, nosuid mount options.\\
\rowcolor{rowalt}
/tmp & \textcolor{cHigh}{recommended} & \textcolor{cOk}{oui} & 1.0 GiB & unbounded & XFS & defaults,nodev,noexec,nosuid,relatime & data & rhel7-security-guide,cis-rhel7 & Separate partition prevents temp files from filling root filesystem. CIS Benchmark requires nodev, nosuid, noexec mount options to restrict execution of binaries and setuid files.\\
/var/tmp & \textcolor{cHigh}{recommended} & \textcolor{cOk}{oui} & 1.0 GiB & unbounded & XFS & defaults,nodev,noexec,nosuid,relatime & data & rhel7-security-guide,cis-rhel7 & Separate partition prevents temp file growth from affecting system stability. CIS Benchmark recommends nodev, noexec, nosuid mount options.\\
\rowcolor{rowalt}
swap & \textcolor{cBlock}{mandatory} & \textcolor{cOk}{oui} & 256 MiB & 128 GiB & swap &  & swap & rhel7-install-guide,rhel7-storage-admin & Automatic partitioning limits swap to 10\% of total disk size or 128 GB maximum. Can be a dedicated swap partition, swap file, or LVM logical volume. For hibernation, size must be at least equal to installed RAM.\\
\bottomrule\end{longtable}}
\subsection{RHEL 8}
\noindent\small \textbf{FS par d\'efaut :} \texttt{XFS} \quad\textbullet\quad \textbf{Swap/paging :} Modern systems with hundreds of GB of RAM should consider swap as a function of system memory workload, not system memory size alone. Minimum 256 MB for systems with 1 GB or less RAM. Swap can be partition, file, or extended LVM logical volume. For hibernation, must be at least equal to RAM size.\normalsize\par\smallskip
\noindent\small\textit{Default filesystem is XFS for / and data partitions. /boot must be ext4 (not XFS). Default partitioning uses LVM. Separate /boot, /home, /tmp, and /var/tmp partitions are recommended for bare-metal installations; optional for virtual machines. Automatic partitioning requires 55 GiB minimum to create /home.}\normalsize\par\smallskip
{\renewcommand{\arraystretch}{1.15}\begin{longtable}{M R S Z Z T O C I N}\toprule\textbf{\normalfont\scriptsize Point de montage / volume} & \textbf{\normalfont\scriptsize Req.} & \textbf{\normalfont\scriptsize Sep.} & \textbf{\normalfont\scriptsize Min} & \textbf{\normalfont\scriptsize Max} & \textbf{\normalfont\scriptsize FS} & \textbf{\normalfont\scriptsize Options de montage} & \textbf{\normalfont\scriptsize Cls} & \textbf{\normalfont\scriptsize Src} & \textbf{\normalfont\scriptsize Notes}\\\midrule\endhead
/ & \textcolor{cBlock}{mandatory} & non & 5.0 GiB & unbounded & XFS & defaults,relatime & system & rhel8-install-guide & Root filesystem. Minimum 5 GiB to accommodate typical installation.\\
\rowcolor{rowalt}
/boot & \textcolor{cBlock}{mandatory} & \textcolor{cOk}{oui} & 1.0 GiB & unbounded & ext4 & defaults,relatime & boot & rhel8-install-guide,rhel8-security-hardening & Must be on physical partition, not LVM. Default size is 1 GB. Recommended 2 GB to avoid running out of space during kernel updates.\\
/boot/efi & \textcolor{cMed}{conditional} & \textcolor{cOk}{oui} & 200 MiB & unbounded & vfat & defaults,umask=0077,shortname=winnt & boot & rhel8-install-guide & Required for UEFI systems on AMD64, Intel 64, and 64-bit ARM. Recommended size 200 MiB. Must not be encrypted.\\
\rowcolor{rowalt}
/home & \textcolor{cHigh}{recommended} & \textcolor{cOk}{oui} & 1.0 GiB & unbounded & XFS & defaults,nodev,nosuid,relatime & data & rhel8-security-hardening,cis-rhel8 & Separate partition protects user data during upgrades and prevents / from filling. Optional for virtual machines.\\
/var & \textcolor{cHigh}{recommended} & \textcolor{cOk}{oui} & 5.0 GiB & unbounded & XFS & defaults,nodev,nosuid,relatime & data & rhel8-install-guide,cis-rhel8 & Holds logs, spool, and transient files. Minimum 5 GiB required. CIS Benchmark recommends nodev and nosuid mount options.\\
\rowcolor{rowalt}
/var/log & \textcolor{cHigh}{recommended} & conditional & 2.0 GiB & unbounded & XFS & defaults,nodev,noexec,nosuid,relatime & data & cis-rhel8,rhel8-security-hardening & Separate partition prevents log growth from affecting system stability. CIS Benchmark recommends nodev, noexec, nosuid mount options.\\
/var/log/audit & \textcolor{cHigh}{recommended} & conditional & 1.0 GiB & unbounded & XFS & defaults,nodev,noexec,nosuid,relatime & data & cis-rhel8,rhel8-security-hardening & Should reside on separate mount point to prevent other processes from consuming audit log space. CIS Benchmark recommends nodev, noexec, nosuid mount options.\\
\rowcolor{rowalt}
/tmp & \textcolor{cHigh}{recommended} & \textcolor{cOk}{oui} & 1.0 GiB & unbounded & XFS & defaults,nodev,noexec,nosuid,relatime & data & rhel8-security-hardening,cis-rhel8 & Separate partition prevents temp files from consuming root filesystem. CIS Benchmark requires nodev, nosuid, noexec mount options.\\
/var/tmp & \textcolor{cHigh}{recommended} & \textcolor{cOk}{oui} & 1.0 GiB & unbounded & XFS & defaults,nodev,noexec,nosuid,relatime & data & rhel8-security-hardening,cis-rhel8 & Separate partition recommended for bare-metal systems. Optional for virtual machines. CIS Benchmark recommends nodev, noexec, nosuid mount options.\\
\rowcolor{rowalt}
swap & \textcolor{cBlock}{mandatory} & \textcolor{cOk}{oui} & 256 MiB & unbounded & swap &  & swap & rhel8-managing-storage & Can be partition, file, or extended LVM logical volume. For hibernation, size must equal or exceed installed RAM. Modern systems should size based on workload, not RAM amount.\\
\bottomrule\end{longtable}}
\subsection{RHEL 9}
\noindent\small \textbf{FS par d\'efaut :} \texttt{XFS} \quad\textbullet\quad \textbf{Swap/paging :} Modern systems with hundreds of GB of RAM should consider swap as a function of system memory workload, not system memory size alone. Minimum 256 MB for systems with 1 GB or less RAM. Swap can be partition, file, or extended LVM logical volume. For hibernation, must be at least equal to RAM size.\normalsize\par\smallskip
\noindent\small\textit{Default filesystem is XFS for / and data partitions. /boot must be ext4 (not XFS). Default partitioning uses LVM. Separate /boot, /home, /tmp, and /var/tmp partitions are recommended for bare-metal installations; optional for virtual machines. Automatic partitioning requires 55 GiB minimum to create /home.}\normalsize\par\smallskip
{\renewcommand{\arraystretch}{1.15}\begin{longtable}{M R S Z Z T O C I N}\toprule\textbf{\normalfont\scriptsize Point de montage / volume} & \textbf{\normalfont\scriptsize Req.} & \textbf{\normalfont\scriptsize Sep.} & \textbf{\normalfont\scriptsize Min} & \textbf{\normalfont\scriptsize Max} & \textbf{\normalfont\scriptsize FS} & \textbf{\normalfont\scriptsize Options de montage} & \textbf{\normalfont\scriptsize Cls} & \textbf{\normalfont\scriptsize Src} & \textbf{\normalfont\scriptsize Notes}\\\midrule\endhead
/ & \textcolor{cBlock}{mandatory} & non & 5.0 GiB & unbounded & XFS & defaults,relatime & system & rhel9-install-guide & Root filesystem. Minimum 5 GiB to accommodate typical installation.\\
\rowcolor{rowalt}
/boot & \textcolor{cBlock}{mandatory} & \textcolor{cOk}{oui} & 1.0 GiB & unbounded & ext4 & defaults,relatime & boot & rhel9-install-guide,rhel9-security-hardening & Must be on physical partition, not LVM. Default size is 1 GB. Must not be encrypted or system cannot boot.\\
/boot/efi & \textcolor{cMed}{conditional} & \textcolor{cOk}{oui} & 200 MiB & unbounded & vfat & defaults,umask=0077,shortname=winnt & boot & rhel9-install-guide & Required for UEFI systems on AMD64, Intel 64, and 64-bit ARM. Recommended size 200 MiB. Must not be encrypted.\\
\rowcolor{rowalt}
/home & \textcolor{cHigh}{recommended} & \textcolor{cOk}{oui} & 1.0 GiB & unbounded & XFS & defaults,nodev,nosuid,relatime & data & rhel9-security-hardening,cis-rhel9 & Separate partition protects user data during upgrades and prevents / from filling. Optional for virtual machines.\\
/var & \textcolor{cHigh}{recommended} & \textcolor{cOk}{oui} & 5.0 GiB & unbounded & XFS & defaults,nodev,nosuid,relatime & data & rhel9-install-guide,cis-rhel9 & Holds logs, spool, and transient files. Minimum 5 GiB required. CIS Benchmark recommends nodev and nosuid mount options.\\
\rowcolor{rowalt}
/var/log & \textcolor{cHigh}{recommended} & conditional & 2.0 GiB & unbounded & XFS & defaults,nodev,noexec,nosuid,relatime & data & cis-rhel9,rhel9-security-hardening & Separate partition prevents log growth from affecting system stability. CIS Benchmark recommends nodev, noexec, nosuid mount options.\\
/var/log/audit & \textcolor{cHigh}{recommended} & conditional & 1.0 GiB & unbounded & XFS & defaults,nodev,noexec,nosuid,relatime & data & cis-rhel9,rhel9-security-hardening & Should reside on separate mount point to prevent other processes from consuming audit log space. CIS Benchmark recommends nodev, noexec, nosuid mount options.\\
\rowcolor{rowalt}
/tmp & \textcolor{cHigh}{recommended} & \textcolor{cOk}{oui} & 1.0 GiB & unbounded & XFS & defaults,nodev,noexec,nosuid,relatime & data & rhel9-security-hardening,cis-rhel9 & Separate partition prevents temp files from consuming root filesystem. CIS Benchmark requires nodev, nosuid, noexec mount options.\\
/var/tmp & \textcolor{cHigh}{recommended} & \textcolor{cOk}{oui} & 1.0 GiB & unbounded & XFS & defaults,nodev,noexec,nosuid,relatime & data & rhel9-security-hardening,cis-rhel9 & Separate partition recommended for bare-metal systems. Optional for virtual machines. CIS Benchmark recommends nodev, noexec, nosuid mount options.\\
\rowcolor{rowalt}
swap & \textcolor{cBlock}{mandatory} & \textcolor{cOk}{oui} & 256 MiB & unbounded & swap &  & swap & rhel9-managing-storage & Can be partition, file, or extended LVM logical volume. For hibernation, size must equal or exceed installed RAM. Modern systems should size based on workload, not RAM amount.\\
\bottomrule\end{longtable}}
\clearpage
\section{Debian/Ubuntu}
\noindent\small\textit{This report covers the five most recent Debian and Ubuntu LTS versions as requested: Debian 11 (Bullseye), Debian 12 (Bookworm), Ubuntu 20.04 LTS, Ubuntu 22.04 LTS, and Ubuntu 24.04 LTS. All versions use ext4 as the default filesystem. The partition recommendations are based on official Debian and Ubuntu installation guides, which show minimal changes across versions. Mount options and separate partition requirements are derived from CIS Benchmarks (the most recent stable versions available for each distribution). The Filesystem Hierarchy Standard (FHS 3.0) provides authoritative guidance on partition purposes and separability. Swap sizing recommendations follow official Ubuntu and Debian documentation, with modern guidance for systems >1GB RAM. CIS Benchmark PDF documents could not be directly retrieved (retrieved=false), but their requirements are well-documented in secondary sources and official documentation. The guidance on zram is specific to Ubuntu 24.04 and later. All partition size recommendations preserve original units from sources (e.g., MiB, GiB, MB, GB). Note: Some CIS Benchmarks are available in multiple versions; the most stable/widely-adopted versions are referenced."}\normalsize\par\medskip
\subsection{Debian 11 (Bullseye)}
\noindent\small \textbf{FS par d\'efaut :} \texttt{ext4} \quad\textbullet\quad \textbf{Swap/paging :} Traditional Debian recommendation: equal to system memory or 1.5-2x RAM. For modern systems (>1GB RAM) without hibernation: minimum round(sqrt(RAM)), maximum 2x RAM. With hibernation: minimum equal to RAM. Minimum 512MB in most cases. Can be swap partition or swapfile.\normalsize\par\smallskip
\noindent\small\textit{Single / partition is acceptable for single-user systems. Multi-user systems should separate /var, /tmp, /home. ext4 is default filesystem for Debian 11. Debian 11 (Bullseye) released August 2021, current LTS version. Official sources show no significant partition layout changes between Debian 11 and 12.}\normalsize\par\smallskip
{\renewcommand{\arraystretch}{1.15}\begin{longtable}{M R S Z Z T O C I N}\toprule\textbf{\normalfont\scriptsize Point de montage / volume} & \textbf{\normalfont\scriptsize Req.} & \textbf{\normalfont\scriptsize Sep.} & \textbf{\normalfont\scriptsize Min} & \textbf{\normalfont\scriptsize Max} & \textbf{\normalfont\scriptsize FS} & \textbf{\normalfont\scriptsize Options de montage} & \textbf{\normalfont\scriptsize Cls} & \textbf{\normalfont\scriptsize Src} & \textbf{\normalfont\scriptsize Notes}\\\midrule\endhead
/ & \textcolor{cBlock}{mandatory} & \textcolor{cOk}{oui} & 8 GiB & unbounded & ext4 & defaults & system & debian-install-guide & Root filesystem containing all essential system files. Minimum 8GB recommended by Debian official documentation.\\
\rowcolor{rowalt}
/boot & \textcolor{cHigh}{recommended} & conditional & 256 MiB & 1 GiB & ext4 & defaults,nosuid,nodev & boot & debian-install-guide,cis-debian-11 & Static boot loader files. Separate partition recommended on GPT/UEFI systems. 256-512MB typical allocation. CIS Benchmark recommends nosuid,nodev mount options.\\
/boot/efi & \textcolor{cMed}{conditional} & \textcolor{cOk}{oui} & 200 MiB & 1 GiB & vfat & defaults & boot & microsoft-uefi-spec & EFI System Partition (ESP) for UEFI firmware. Required for UEFI boot mode only. Minimum 200MB per UEFI specification (Microsoft Learn), typically 512MB on Linux systems.\\
\rowcolor{rowalt}
/home & \textcolor{cHigh}{recommended} & \textcolor{cOk}{oui} & unbounded & unbounded & ext4 & defaults,nosuid,nodev & data & debian-install-guide,cis-debian-11 & User home directories. Recommended separate partition for multi-user systems or large disk systems. CIS Benchmark recommends nosuid,nodev mount options.\\
/usr & \textcolor{cLow}{optional} & conditional & unbounded & unbounded & ext4 & defaults & system & fhs-3.0 & Static user programs and read-only data. Optional separate partition. Can be mounted read-only if separate filesystem (FHS 3.0 design principle).\\
\rowcolor{rowalt}
/var & \textcolor{cHigh}{recommended} & \textcolor{cOk}{oui} & 200 MiB & unbounded & ext4 & defaults,nosuid,nodev & data & debian-install-guide,cis-debian-11 & Variable data: logs, APT cache, dpkg information. Recommended minimum 200MB for base system. Recommended separate for multi-user systems. CIS Benchmark recommends nosuid,nodev.\\
/var/log & \textcolor{cHigh}{recommended} & conditional & 100 MiB & unbounded & ext4 & defaults,nodev,nosuid,noexec & data & cis-debian-11 & System log files. CIS Benchmark recommends separate partition with nodev,nosuid,noexec mount options to prevent log filling and executable code execution.\\
\rowcolor{rowalt}
/var/tmp & \textcolor{cHigh}{recommended} & conditional & 100 MiB & unbounded & ext4 & defaults,nodev,nosuid,noexec & data & cis-debian-11 & Temporary files preserved between reboots. CIS Benchmark recommends separate partition with nodev,nosuid,noexec mount options.\\
/tmp & \textcolor{cHigh}{recommended} & \textcolor{cOk}{oui} & 500 MiB & 2 GiB & ext4 & defaults,nodev,nosuid,noexec & data & cis-debian-11 & Temporary files deleted on reboot. CIS Benchmark requires separate partition with nodev,nosuid,noexec mount options to prevent executable code and protect against resource exhaustion.\\
\rowcolor{rowalt}
/dev/shm & \textcolor{cLow}{optional} & non & N/A & N/A & tmpfs & defaults,nodev,nosuid,noexec & data & cis-debian-11 & Shared memory filesystem (tmpfs). CIS Benchmark recommends mounting with nodev,nosuid,noexec options to prevent malicious executable placement.\\
swap & \textcolor{cBlock}{mandatory} & \textcolor{cOk}{oui} & 512 MiB & 2x RAM & swap & defaults & swap & debian-install-guide & Swap space for virtual memory. Traditional guideline: equal to system RAM. Minimum 512MB in most cases. Without hibernation: min=round(sqrt(RAM)), max=2xRAM.\\
\bottomrule\end{longtable}}
\subsection{Debian 12 (Bookworm)}
\noindent\small \textbf{FS par d\'efaut :} \texttt{ext4} \quad\textbullet\quad \textbf{Swap/paging :} Traditional Debian recommendation: equal to system memory or 1.5-2x RAM. For modern systems (>1GB RAM) without hibernation: minimum round(sqrt(RAM)), maximum 2x RAM. With hibernation: minimum equal to RAM. Minimum 512MB in most cases. Can be swap partition or swapfile.\normalsize\par\smallskip
\noindent\small\textit{Debian 12 (Bookworm) released June 2023. No significant changes from Debian 11 in recommended partition layout. ext4 remains default filesystem. CIS Benchmark v1.1.0 available for Debian 12.}\normalsize\par\smallskip
{\renewcommand{\arraystretch}{1.15}\begin{longtable}{M R S Z Z T O C I N}\toprule\textbf{\normalfont\scriptsize Point de montage / volume} & \textbf{\normalfont\scriptsize Req.} & \textbf{\normalfont\scriptsize Sep.} & \textbf{\normalfont\scriptsize Min} & \textbf{\normalfont\scriptsize Max} & \textbf{\normalfont\scriptsize FS} & \textbf{\normalfont\scriptsize Options de montage} & \textbf{\normalfont\scriptsize Cls} & \textbf{\normalfont\scriptsize Src} & \textbf{\normalfont\scriptsize Notes}\\\midrule\endhead
/ & \textcolor{cBlock}{mandatory} & \textcolor{cOk}{oui} & 8 GiB & unbounded & ext4 & defaults & system & debian-install-guide & Root filesystem containing all essential system files. Minimum 8GB recommended by Debian official documentation.\\
\rowcolor{rowalt}
/boot & \textcolor{cHigh}{recommended} & conditional & 256 MiB & 1 GiB & ext4 & defaults,nosuid,nodev & boot & debian-install-guide,cis-debian-12 & Static boot loader files. Separate partition recommended on GPT/UEFI systems. 256-512MB typical allocation. CIS Benchmark recommends nosuid,nodev mount options.\\
/boot/efi & \textcolor{cMed}{conditional} & \textcolor{cOk}{oui} & 200 MiB & 1 GiB & vfat & defaults & boot & microsoft-uefi-spec & EFI System Partition (ESP) for UEFI firmware. Required for UEFI boot mode only. Minimum 200MB per UEFI specification, typically 512MB on Linux systems.\\
\rowcolor{rowalt}
/home & \textcolor{cHigh}{recommended} & \textcolor{cOk}{oui} & unbounded & unbounded & ext4 & defaults,nosuid,nodev & data & debian-install-guide,cis-debian-12 & User home directories. Recommended separate partition for multi-user systems or large disk systems. CIS Benchmark recommends nosuid,nodev mount options.\\
/usr & \textcolor{cLow}{optional} & conditional & unbounded & unbounded & ext4 & defaults & system & fhs-3.0 & Static user programs and read-only data. Optional separate partition. Can be mounted read-only if separate filesystem (FHS 3.0 design principle).\\
\rowcolor{rowalt}
/var & \textcolor{cHigh}{recommended} & \textcolor{cOk}{oui} & 200 MiB & unbounded & ext4 & defaults,nosuid,nodev & data & debian-install-guide,cis-debian-12 & Variable data: logs, APT cache, dpkg information. Recommended minimum 200MB for base system. Recommended separate for multi-user systems. CIS Benchmark recommends nosuid,nodev.\\
/var/log & \textcolor{cHigh}{recommended} & conditional & 100 MiB & unbounded & ext4 & defaults,nodev,nosuid,noexec & data & cis-debian-12 & System log files. CIS Benchmark recommends separate partition with nodev,nosuid,noexec mount options to prevent log filling and executable code execution.\\
\rowcolor{rowalt}
/var/tmp & \textcolor{cHigh}{recommended} & conditional & 100 MiB & unbounded & ext4 & defaults,nodev,nosuid,noexec & data & cis-debian-12 & Temporary files preserved between reboots. CIS Benchmark recommends separate partition with nodev,nosuid,noexec mount options.\\
/tmp & \textcolor{cHigh}{recommended} & \textcolor{cOk}{oui} & 500 MiB & 2 GiB & ext4 & defaults,nodev,nosuid,noexec & data & cis-debian-12 & Temporary files deleted on reboot. CIS Benchmark requires separate partition with nodev,nosuid,noexec mount options to prevent executable code and protect against resource exhaustion.\\
\rowcolor{rowalt}
/dev/shm & \textcolor{cLow}{optional} & non & N/A & N/A & tmpfs & defaults,nodev,nosuid,noexec & data & cis-debian-12 & Shared memory filesystem (tmpfs). CIS Benchmark recommends mounting with nodev,nosuid,noexec options to prevent malicious executable placement.\\
swap & \textcolor{cBlock}{mandatory} & \textcolor{cOk}{oui} & 512 MiB & 2x RAM & swap & defaults & swap & debian-install-guide & Swap space for virtual memory. Traditional guideline: equal to system RAM. Minimum 512MB in most cases. Without hibernation: min=round(sqrt(RAM)), max=2xRAM.\\
\bottomrule\end{longtable}}
\subsection{Ubuntu 20.04 LTS (Focal)}
\noindent\small \textbf{FS par d\'efaut :} \texttt{ext4} \quad\textbullet\quad \textbf{Swap/paging :} Ubuntu official recommendation: For systems with hibernation support, swap partition must be at least as large as RAM. For systems without hibernation: minimum round(sqrt(RAM)), maximum 2x RAM. Historically: equal to system memory. Modern guidance emphasizes that excessive swap is inefficient since disk access is \textasciitilde{}1000x slower than RAM.\normalsize\par\smallskip
\noindent\small\textit{Ubuntu 20.04 LTS (Focal Fossa) released April 2020, standard LTS with 5-year support. CIS Benchmark v1.1.0 available. Single / partition with /home separate is acceptable for most users.}\normalsize\par\smallskip
{\renewcommand{\arraystretch}{1.15}\begin{longtable}{M R S Z Z T O C I N}\toprule\textbf{\normalfont\scriptsize Point de montage / volume} & \textbf{\normalfont\scriptsize Req.} & \textbf{\normalfont\scriptsize Sep.} & \textbf{\normalfont\scriptsize Min} & \textbf{\normalfont\scriptsize Max} & \textbf{\normalfont\scriptsize FS} & \textbf{\normalfont\scriptsize Options de montage} & \textbf{\normalfont\scriptsize Cls} & \textbf{\normalfont\scriptsize Src} & \textbf{\normalfont\scriptsize Notes}\\\midrule\endhead
/ & \textcolor{cBlock}{mandatory} & \textcolor{cOk}{oui} & 8 GiB & unbounded & ext4 & defaults & system & ubuntu-disk-space & Root filesystem. Minimum 8GB, recommended 15GB by Ubuntu official documentation.\\
\rowcolor{rowalt}
/boot & \textcolor{cHigh}{recommended} & conditional & 256 MiB & 1 GiB & ext4 & defaults,nosuid,nodev & boot & ubuntu-partitioning & Boot loader files. Separate partition recommended on GPT/UEFI systems. Can be necessary if /boot files are >100GB from disk start (some BIOS limitations).\\
/boot/efi & \textcolor{cMed}{conditional} & \textcolor{cOk}{oui} & 200 MiB & 1 GiB & vfat & defaults & boot & ubuntu-partitioning & EFI System Partition (ESP) created automatically by installer in UEFI mode. Ubuntu 20.04 creates minimum 538MiB on UEFI systems.\\
\rowcolor{rowalt}
/home & \textcolor{cHigh}{recommended} & \textcolor{cOk}{oui} & unbounded & unbounded & ext4 & defaults,nosuid,nodev & data & ubuntu-partitioning,cis-ubuntu-20 & User home directories. Recommended for multi-user systems. CIS Benchmark for Ubuntu 20.04 recommends nosuid,nodev mount options.\\
/usr & \textcolor{cLow}{optional} & conditional & 10 GiB & unbounded & ext4 & defaults & system & ubuntu-disk-space & User programs and data. Ubuntu documentation recommends \textasciitilde{}10GB for separate /usr partition. Can be mounted read-only if separate.\\
\rowcolor{rowalt}
/var & \textcolor{cHigh}{recommended} & \textcolor{cOk}{oui} & 2 GiB & unbounded & ext4 & defaults,nosuid,nodev & data & ubuntu-disk-space,cis-ubuntu-20 & Variable data. Ubuntu recommends 2GB for separate /var partition, 3GB+ if upgrading packages. CIS Benchmark recommends nosuid,nodev.\\
/var/log & \textcolor{cHigh}{recommended} & conditional & 100 MiB & unbounded & ext4 & defaults,nodev,nosuid,noexec & data & cis-ubuntu-20 & System log files. CIS Benchmark for Ubuntu 20.04 requires separate partition with nodev,nosuid,noexec mount options.\\
\rowcolor{rowalt}
/var/tmp & \textcolor{cHigh}{recommended} & conditional & 100 MiB & unbounded & ext4 & defaults,nodev,nosuid,noexec & data & cis-ubuntu-20 & Temporary files preserved between reboots. CIS Benchmark requires separate partition with nodev,nosuid,noexec mount options.\\
/tmp & \textcolor{cHigh}{recommended} & \textcolor{cOk}{oui} & 500 MiB & 2 GiB & ext4 & defaults,nodev,nosuid,noexec & data & cis-ubuntu-20 & Temporary files. Ubuntu recommends 20-50MB minimum historically, but CIS Benchmark requires separate partition with nodev,nosuid,noexec mount options.\\
\rowcolor{rowalt}
/dev/shm & \textcolor{cLow}{optional} & non & N/A & N/A & tmpfs & defaults,nodev,nosuid,noexec & data & cis-ubuntu-20 & Shared memory filesystem (tmpfs). CIS Benchmark recommends mounting with nodev,nosuid,noexec options.\\
swap & \textcolor{cBlock}{mandatory} & \textcolor{cOk}{oui} & 1 GiB & 2x RAM & swap & defaults & swap & ubuntu-swap-faq & Swap space. Ubuntu official guidance: equal to system RAM. For modern systems (>1GB RAM) without hibernation: min=round(sqrt(RAM)), max=2xRAM. With hibernation: min=RAM size.\\
\bottomrule\end{longtable}}
\subsection{Ubuntu 22.04 LTS (Jammy)}
\noindent\small \textbf{FS par d\'efaut :} \texttt{ext4} \quad\textbullet\quad \textbf{Swap/paging :} Ubuntu 22.04 LTS recommends: For systems with hibernation, swap partition must be at least as large as RAM. For systems without hibernation: minimum round(sqrt(RAM)), maximum 2x RAM. Modern systems with abundant RAM can use less swap. Swapfile is viable alternative to swap partition.\normalsize\par\smallskip
\noindent\small\textit{Ubuntu 22.04 LTS (Jammy Jellyfish) released April 2022, standard LTS with 5-year support (until 2027). No major partition layout changes from 20.04. ext4 remains default. CIS Benchmark v1.0.0 and later versions available.}\normalsize\par\smallskip
{\renewcommand{\arraystretch}{1.15}\begin{longtable}{M R S Z Z T O C I N}\toprule\textbf{\normalfont\scriptsize Point de montage / volume} & \textbf{\normalfont\scriptsize Req.} & \textbf{\normalfont\scriptsize Sep.} & \textbf{\normalfont\scriptsize Min} & \textbf{\normalfont\scriptsize Max} & \textbf{\normalfont\scriptsize FS} & \textbf{\normalfont\scriptsize Options de montage} & \textbf{\normalfont\scriptsize Cls} & \textbf{\normalfont\scriptsize Src} & \textbf{\normalfont\scriptsize Notes}\\\midrule\endhead
/ & \textcolor{cBlock}{mandatory} & \textcolor{cOk}{oui} & 8 GiB & unbounded & ext4 & defaults & system & ubuntu-disk-space & Root filesystem. Minimum 8GB, recommended 15GB by Ubuntu official documentation.\\
\rowcolor{rowalt}
/boot & \textcolor{cHigh}{recommended} & conditional & 256 MiB & 1 GiB & ext4 & defaults,nosuid,nodev & boot & ubuntu-partitioning & Boot loader files. Separate partition recommended on GPT/UEFI systems. Similar to Ubuntu 20.04.\\
/boot/efi & \textcolor{cMed}{conditional} & \textcolor{cOk}{oui} & 200 MiB & 1 GiB & vfat & defaults & boot & ubuntu-partitioning & EFI System Partition (ESP) for UEFI systems. Created automatically by installer.\\
\rowcolor{rowalt}
/home & \textcolor{cHigh}{recommended} & \textcolor{cOk}{oui} & unbounded & unbounded & ext4 & defaults,nosuid,nodev & data & ubuntu-partitioning,cis-ubuntu-22 & User home directories. Recommended for multi-user systems. CIS Benchmark for Ubuntu 22.04 recommends nosuid,nodev mount options.\\
/usr & \textcolor{cLow}{optional} & conditional & 10 GiB & unbounded & ext4 & defaults & system & ubuntu-disk-space & User programs and data. Can be mounted read-only if separate filesystem.\\
\rowcolor{rowalt}
/var & \textcolor{cHigh}{recommended} & \textcolor{cOk}{oui} & 2 GiB & unbounded & ext4 & defaults,nosuid,nodev & data & ubuntu-disk-space,cis-ubuntu-22 & Variable data. Ubuntu recommends 2GB minimum for separate /var. CIS Benchmark recommends nosuid,nodev.\\
/var/log & \textcolor{cHigh}{recommended} & conditional & 100 MiB & unbounded & ext4 & defaults,nodev,nosuid,noexec & data & cis-ubuntu-22 & System log files. CIS Benchmark for Ubuntu 22.04 requires separate partition with nodev,nosuid,noexec mount options.\\
\rowcolor{rowalt}
/var/tmp & \textcolor{cHigh}{recommended} & conditional & 100 MiB & unbounded & ext4 & defaults,nodev,nosuid,noexec & data & cis-ubuntu-22 & Temporary files preserved between reboots. CIS Benchmark requires separate partition with nodev,nosuid,noexec mount options.\\
/tmp & \textcolor{cHigh}{recommended} & \textcolor{cOk}{oui} & 500 MiB & 2 GiB & ext4 & defaults,nodev,nosuid,noexec & data & cis-ubuntu-22 & Temporary files. CIS Benchmark requires separate partition with nodev,nosuid,noexec mount options.\\
\rowcolor{rowalt}
/dev/shm & \textcolor{cLow}{optional} & non & N/A & N/A & tmpfs & defaults,nodev,nosuid,noexec & data & cis-ubuntu-22 & Shared memory filesystem (tmpfs). CIS Benchmark recommends mounting with nodev,nosuid,noexec options.\\
swap & \textcolor{cBlock}{mandatory} & \textcolor{cOk}{oui} & 1 GiB & 2x RAM & swap & defaults & swap & ubuntu-swap-faq & Swap space. Ubuntu official guidance: For modern systems (>1GB RAM) without hibernation: min=round(sqrt(RAM)), max=2xRAM. With hibernation: min=RAM size. Can also use swapfile.\\
\bottomrule\end{longtable}}
\subsection{Ubuntu 24.04 LTS (Noble)}
\noindent\small \textbf{FS par d\'efaut :} \texttt{ext4} \quad\textbullet\quad \textbf{Swap/paging :} Ubuntu 24.04 LTS: For systems with hibernation, swap must be at least RAM size. For systems without hibernation: min=round(sqrt(RAM)), max=2xRAM. Ubuntu 24.04 supports zram (compressed in-memory swap) via zram-tools package, which can compress data 2-3x. zram prevents hibernation; use zswap if sleep-to-disk is needed. Modern systems with >8GB RAM can use smaller swap or zram only.\normalsize\par\smallskip
\noindent\small\textit{Ubuntu 24.04 LTS (Noble Numbat) released April 2024, latest LTS with 5-year standard support (until 2029). Introduces zram support for memory-constrained systems. CIS Benchmark v1.0.0 and v2.0.0 available. Can use LVM for flexible storage management.}\normalsize\par\smallskip
{\renewcommand{\arraystretch}{1.15}\begin{longtable}{M R S Z Z T O C I N}\toprule\textbf{\normalfont\scriptsize Point de montage / volume} & \textbf{\normalfont\scriptsize Req.} & \textbf{\normalfont\scriptsize Sep.} & \textbf{\normalfont\scriptsize Min} & \textbf{\normalfont\scriptsize Max} & \textbf{\normalfont\scriptsize FS} & \textbf{\normalfont\scriptsize Options de montage} & \textbf{\normalfont\scriptsize Cls} & \textbf{\normalfont\scriptsize Src} & \textbf{\normalfont\scriptsize Notes}\\\midrule\endhead
/ & \textcolor{cBlock}{mandatory} & \textcolor{cOk}{oui} & 8 GiB & unbounded & ext4 & defaults & system & ubuntu-system-requirements & Root filesystem. Minimum 8GB, recommended 15GB. Ubuntu 24.04 Server requires minimum 5GB (ISO install), Desktop requires 25GB.\\
\rowcolor{rowalt}
/boot & \textcolor{cHigh}{recommended} & conditional & 256 MiB & 1 GiB & ext4 & defaults,nosuid,nodev & boot & ubuntu-partitioning & Boot loader files. Separate partition recommended on GPT/UEFI systems.\\
/boot/efi & \textcolor{cMed}{conditional} & \textcolor{cOk}{oui} & 200 MiB & 1 GiB & vfat & defaults & boot & ubuntu-partitioning & EFI System Partition (ESP) for UEFI systems. Created automatically by installer in UEFI mode.\\
\rowcolor{rowalt}
/home & \textcolor{cHigh}{recommended} & \textcolor{cOk}{oui} & unbounded & unbounded & ext4 & defaults,nosuid,nodev & data & ubuntu-partitioning,cis-ubuntu-24 & User home directories. Recommended for multi-user systems. CIS Benchmark for Ubuntu 24.04 recommends nosuid,nodev mount options.\\
/usr & \textcolor{cLow}{optional} & conditional & 10 GiB & unbounded & ext4 & defaults & system & ubuntu-disk-space & User programs and data. Can be mounted read-only if separate filesystem.\\
\rowcolor{rowalt}
/var & \textcolor{cHigh}{recommended} & \textcolor{cOk}{oui} & 2 GiB & unbounded & ext4 & defaults,nosuid,nodev & data & ubuntu-disk-space,cis-ubuntu-24 & Variable data. Ubuntu recommends 2GB minimum for separate /var. CIS Benchmark recommends nosuid,nodev.\\
/var/log & \textcolor{cHigh}{recommended} & conditional & 100 MiB & unbounded & ext4 & defaults,nodev,nosuid,noexec & data & cis-ubuntu-24 & System log files. CIS Benchmark for Ubuntu 24.04 requires separate partition with nodev,nosuid,noexec mount options.\\
\rowcolor{rowalt}
/var/tmp & \textcolor{cHigh}{recommended} & conditional & 100 MiB & unbounded & ext4 & defaults,nodev,nosuid,noexec & data & cis-ubuntu-24 & Temporary files preserved between reboots. CIS Benchmark requires separate partition with nodev,nosuid,noexec mount options.\\
/tmp & \textcolor{cHigh}{recommended} & \textcolor{cOk}{oui} & 500 MiB & 2 GiB & ext4 & defaults,nodev,nosuid,noexec & data & cis-ubuntu-24 & Temporary files. CIS Benchmark requires separate partition with nodev,nosuid,noexec mount options.\\
\rowcolor{rowalt}
/dev/shm & \textcolor{cLow}{optional} & non & N/A & N/A & tmpfs & defaults,nodev,nosuid,noexec & data & cis-ubuntu-24 & Shared memory filesystem (tmpfs). CIS Benchmark recommends mounting with nodev,nosuid,noexec options.\\
swap & \textcolor{cBlock}{mandatory} & \textcolor{cOk}{oui} & 1 GiB & 2x RAM & swap & defaults & swap & ubuntu-swap-faq & Swap space or zram. Ubuntu 24.04 supports zram for compressed in-memory swap. Traditional swap partition or swapfile also supported.\\
\bottomrule\end{longtable}}
\clearpage
\section{SUSE/SLES}
\noindent\small\textit{Complete coverage of SLES 12 SP5 and SLES 15 (SP1-SP7) official storage partitioning from SUSE documentation. Default Btrfs configuration with subvolumes documented from official deployment, storage admin, and filesystem guides. CIS Benchmark partition and mount option requirements sourced from CIS Security website. FHS 3.0 principles documented from Linux Foundation. Note: CIS SUSE Benchmark PDFs were not fully readable in binary format; partition/mount recommendations inferred from general CIS Linux principles and verified against SLES-specific documentation where available."}\normalsize\par\medskip
\subsection{SLES 12 SP5}
\noindent\small \textbf{FS par d\'efaut :} \texttt{Btrfs (root); XFS (additional partitions)} \quad\textbullet\quad \textbf{Swap/paging :} Minimum 256 MiB. For suspend-to-disk/hibernation, allocate 512 MB-1 GB. Modern guidance: for RAM > 2GB, 20-25\% of RAM is acceptable; systems with 4GB+ RAM can use proportionally smaller swap. If swap fills up, consider adding system RAM instead.\normalsize\par\smallskip
\noindent\small\textit{SLES 12 SP5 uses Btrfs with comprehensive subvolume layout by default. Root is @ subvolume. Key subvolumes excluded from snapshots: @/home, @/var, @/opt, @/srv, @/usr/local, @/tmp, @/boot/grub2/*, @/var/log (NoCOW). CoW disabled on /var/log and database directories (@/var/lib/mariadb, @/var/lib/mysql, @/var/lib/pgqsl, @/var/lib/libvirt/images) for performance. XFS is default for non-root partitions if separated. System rollbacks supported only if default subvolume configuration not modified.}\normalsize\par\smallskip
{\renewcommand{\arraystretch}{1.15}\begin{longtable}{M R S Z Z T O C I N}\toprule\textbf{\normalfont\scriptsize Point de montage / volume} & \textbf{\normalfont\scriptsize Req.} & \textbf{\normalfont\scriptsize Sep.} & \textbf{\normalfont\scriptsize Min} & \textbf{\normalfont\scriptsize Max} & \textbf{\normalfont\scriptsize FS} & \textbf{\normalfont\scriptsize Options de montage} & \textbf{\normalfont\scriptsize Cls} & \textbf{\normalfont\scriptsize Src} & \textbf{\normalfont\scriptsize Notes}\\\midrule\endhead
/ & \textcolor{cBlock}{mandatory} & non & 16 GiB & unbounded & Btrfs & defaults,relatime & system & suse-sles12-sp5-deployment & Btrfs with subvolumes (@, @/var, @/home, @/opt, @/srv, @/usr/local, @/tmp, @/boot/grub2/*, @/.snapshots). Root subvolume is @. Default size 16 GiB minimum (with snapshots), recommended 32 GiB or more if /home is not separate. Snapshots enabled via Snapper infrastructure.\\
\rowcolor{rowalt}
/boot & \textcolor{cHigh}{recommended} & conditional & 500 MiB & 1 GiB & Ext4 & defaults,relatime & boot & suse-sles12-sp5-expert-partitioner & 500 MB recommended for most deployments. /boot may be included in root Btrfs by default but can be separated for additional security. Ext4 is recommended for boot partitions.\\
/boot/efi & \textcolor{cMed}{conditional} & \textcolor{cOk}{oui} & 260 MiB & 500 MiB & FAT32 (VFAT) & defaults,umask=0077 & boot & suse-sles12-sp5-uefi & UEFI/EFI systems only. Must be formatted with FAT32. Can be RAID1 for redundancy. 260 MiB minimum ensures proper FAT32 formatting. 500 MiB common guideline.\\
\rowcolor{rowalt}
/var & \textcolor{cHigh}{recommended} & conditional & 5 GiB & unbounded & Btrfs subvolume or XFS & defaults,nodev,nosuid,relatime & data & suse-sles12-sp5-filesystems & In SLES 12, /var may be a Btrfs subvolume (@/var) with CoW disabled (NoCOW attribute set) to improve performance. /var/log has NoCOW set by default. If separate partition, XFS is default. CIS Benchmark recommends separate partition with nodev, nosuid.\\
/tmp & \textcolor{cHigh}{recommended} & conditional & 1 GiB & unbounded & Ext4 or XFS & defaults,nodev,nosuid,noexec,relatime & data & cis-suse-benchmark & CIS Benchmark 1.1.2-1.1.5: Recommend separate partition with nodev, nosuid, noexec options to prevent code execution. /tmp subvolume is excluded from snapshots.\\
\rowcolor{rowalt}
/home & \textcolor{cHigh}{recommended} & conditional & 10 GiB & unbounded & XFS (if separate) & defaults,nodev,nosuid,relatime & data & suse-sles12-sp5-filesystems & CIS Benchmark recommends separate partition. If on Btrfs (@/home subvolume), excluded from snapshots to prevent data loss. Default may be separate XFS partition or Btrfs subvolume.\\
/var/tmp & \textcolor{cLow}{optional} & conditional & 1 GiB & unbounded & Ext4 or XFS & defaults,nodev,nosuid,noexec,relatime & data & cis-suse-benchmark & CIS Benchmark recommends separate partition with nodev, nosuid, noexec. Often point to /tmp via symbolic link. May be Btrfs subvolume (@/var/tmp) with CoW.\\
\rowcolor{rowalt}
/dev/shm & \textcolor{cHigh}{recommended} & \textcolor{cOk}{oui} & 512 MiB & unbounded & tmpfs & defaults,nodev,nosuid,noexec,relatime,size=<percent-of-RAM> & system & cis-suse-benchmark & CIS Benchmark 1.1.6: Recommend mount options nodev, nosuid, noexec. Tmpfs in-memory filesystem; size can be specified as percentage of RAM.\\
swap & \textcolor{cHigh}{recommended} & \textcolor{cOk}{oui} & 256 MiB & unbounded & swap & defaults,pri=-2 & swap & suse-sles12-sp5-expert-partitioner & Minimum 256 MiB. For suspend-to-disk, 512 MB-1 GB. Modern systems with adequate RAM may prioritize memory over swap. If RAM > 2GB, can be smaller than historical 2x RAM guideline.\\
\bottomrule\end{longtable}}
\subsection{SLES 15 (SP1-SP7)}
\noindent\small \textbf{FS par d\'efaut :} \texttt{Btrfs (root); XFS (additional partitions)} \quad\textbullet\quad \textbf{Swap/paging :} Minimum 256 MiB. For suspend-to-disk/hibernation, allocate 512 MB-1 GB. Modern guidance (SLES 15): for systems with adequate RAM, swap can be smaller than historical 2x RAM. Consider 20-25\% of RAM for systems with 4GB+ RAM. Systems with more than adequate RAM may not need swap at all.\normalsize\par\smallskip
\noindent\small\textit{SLES 15 (all SP versions) uses Btrfs with refined subvolume layout. Root is @ subvolume. Key change from SLES 12: /var is now a dedicated Btrfs subvolume (@/var) with Copy-On-Write DISABLED for performance (logs, caches, VM images, databases). Subvolumes excluded from snapshots: @/var (entire), @/home, @/opt, @/srv, @/usr/local, @/tmp, @/boot/grub2/* (architecture-specific), @/.snapshots. XFS is default for additional non-root partitions if separated. Snapper snapshots enabled for root filesystem. System rollbacks only supported if default subvolume configuration maintained. FHS 3.0 principles followed for /tmp, /var, /home separability.}\normalsize\par\smallskip
{\renewcommand{\arraystretch}{1.15}\begin{longtable}{M R S Z Z T O C I N}\toprule\textbf{\normalfont\scriptsize Point de montage / volume} & \textbf{\normalfont\scriptsize Req.} & \textbf{\normalfont\scriptsize Sep.} & \textbf{\normalfont\scriptsize Min} & \textbf{\normalfont\scriptsize Max} & \textbf{\normalfont\scriptsize FS} & \textbf{\normalfont\scriptsize Options de montage} & \textbf{\normalfont\scriptsize Cls} & \textbf{\normalfont\scriptsize Src} & \textbf{\normalfont\scriptsize Notes}\\\midrule\endhead
/ & \textcolor{cBlock}{mandatory} & non & 16 GiB & unbounded & Btrfs & defaults,relatime & system & suse-sles15-sp7-expert-partitioner & Btrfs with subvolumes (@, @/var, @/home, @/opt, @/srv, @/usr/local, @/tmp, @/boot/grub2/*, @/.snapshots). Root subvolume is @. Minimum 16 GiB (with snapshots), SUSE recommends 32 GiB or more if /home not separate. Snapshots via Snapper. In SLES 15, /var is dedicated subvolume with CoW disabled.\\
\rowcolor{rowalt}
/boot & \textcolor{cHigh}{recommended} & conditional & 500 MiB & 1 GiB & Ext4 or Btrfs & defaults,relatime & boot & suse-sles15-sp7-expert-partitioner & 500 MB recommended. Can be separate partition or included in root. Ext4 recommended for boot if separate. Btrfs boot support also available.\\
/boot/efi & \textcolor{cMed}{conditional} & \textcolor{cOk}{oui} & 260 MiB & 500 MiB & FAT32 (VFAT) & defaults,umask=0077 & boot & suse-sles15-sp7-uefi & UEFI systems only. Must be FAT32 formatted. Physical partition or RAID1 only (no LVM, no other RAID). 260 MiB minimum for FAT32 compatibility, 500 MiB guideline.\\
\rowcolor{rowalt}
/var & \textcolor{cHigh}{recommended} & conditional & 5 GiB & unbounded & Btrfs subvolume or XFS & defaults,nodev,nosuid,relatime & data & suse-sles15-sp7-filesystems & In SLES 15, /var is dedicated Btrfs subvolume (@/var) with Copy-On-Write disabled (contains logs, caches, VM images, databases). Excluded from snapshots. If separate partition, XFS default. CIS Benchmark recommends separate with nodev, nosuid. Do NOT use noexec on /var.\\
/tmp & \textcolor{cHigh}{recommended} & conditional & 1 GiB & unbounded & Ext4, XFS, or tmpfs & defaults,nodev,nosuid,noexec,relatime & data & cis-suse-benchmark & CIS Benchmark 1.1.2-1.1.5: Recommend separate partition with nodev, nosuid, noexec. Btrfs @/tmp subvolume excluded from snapshots. If tmpfs: defaults,rw,nosuid,nodev,noexec,relatime,size=<X>.\\
\rowcolor{rowalt}
/home & \textcolor{cHigh}{recommended} & conditional & 10 GiB & unbounded & XFS (if separate) or Btrfs & defaults,nodev,nosuid,relatime & data & suse-sles15-sp7-filesystems & CIS Benchmark recommends separate partition. If Btrfs subvolume (@/home), excluded from snapshots to prevent user data loss. Default may be separate XFS or Btrfs subvolume.\\
/var/tmp & \textcolor{cLow}{optional} & conditional & 1 GiB & unbounded & Ext4, XFS, or Btrfs & defaults,nodev,nosuid,noexec,relatime & data & cis-suse-benchmark & CIS Benchmark: separate partition with nodev, nosuid, noexec. Often symlink to /tmp. May be Btrfs subvolume.\\
\rowcolor{rowalt}
/dev/shm & \textcolor{cHigh}{recommended} & \textcolor{cOk}{oui} & 512 MiB & unbounded & tmpfs & defaults,nodev,nosuid,noexec,relatime,size=<percent-of-RAM> & system & cis-suse-benchmark & CIS Benchmark 1.1.6: Mount options nodev, nosuid, noexec. Tmpfs in-memory filesystem.\\
swap & \textcolor{cHigh}{recommended} & \textcolor{cOk}{oui} & 256 MiB & unbounded & swap & defaults,pri=-2 & swap & suse-sles15-sp7-expert-partitioner & Minimum 256 MiB. For suspend-to-disk, 512 MB-1 GB. Modern systems with adequate RAM can use proportionally smaller swap. If swap fills, add RAM instead.\\
\bottomrule\end{longtable}}
\clearpage
\section{AIX}
\noindent\small\textit{This reference covers the official IBM AIX storage partitioning standards for versions 7.1, 7.2, and 7.3. Partition data is derived from IBM official documentation (release notes, system administration guides) and IBM support pages. All twelve rootvg logical volumes are documented: hd1 (/home), hd2 (/usr), hd3 (/tmp), hd4 (/), hd5 (boot), hd6 (paging/swap), hd8 (JFS2 log), hd9var (/var), hd10opt (/opt), hd11admin (/admin), livedump (/var/adm/ras/livedump), and lg\_dumplv (dump). Key version-specific differences noted: AIX 7.3 increased hd5 (boot) minimum from 24 MB to 40 MB; all versions default to 512 MB hd6 paging space and 20 GB total system disk. JFS2 is the universal filesystem type. Mount options default to 'rw' with no hardening options documented in standard installation; CIS Benchmark reference available but legacy (pre-AIX 7.x version). Paging space sizing recommendations provided (2x RAM up to 256 MB, then 512 MB + (RAM-256MB)x1.25). Datavg (separate data volume group) is established best practice for application data separation. Retrieved sources are confirmed accessible; some IBM docs pages return 403 Forbidden, marked as retrieved=false. This is a read-only research report compiled from publicly accessible IBM documentation."
}\normalsize\par\medskip
\subsection{7.1}
\noindent\small \textbf{FS par d\'efaut :} \texttt{JFS2} \quad\textbullet\quad \textbf{Swap/paging :} The general recommendation is that the sum of the sizes of the paging spaces should be equal to at least twice the size of real memory, up to 256 MB of memory (512 MB paging). For systems with memory > 256 MB: total paging space = 512 MB + (memory size - 256 MB) * 1.25. Each single paging space cannot exceed 64 GB. Default hd6 is 512 MB.\normalsize\par\smallskip
\noindent\small\textit{AIX 7.1 requires minimum 10 GB disk for default installation. All standard filesystems must be JFS2. The rootvg contains mandatory system filesystems (hd2, hd3, hd4, hd5, hd6, hd8, hd9var). Optional filesystems include hd1, hd10opt, hd11admin. All rootvg filesystems except hd10opt cannot be moved. Default paging space (hd6) of 512 MB often insufficient for production systems. Application data should reside in separate datavg volume group. AIX 7.1 baseline documentation available from IBM.}\normalsize\par\smallskip
{\renewcommand{\arraystretch}{1.15}\begin{longtable}{M R S Z Z T O C I N}\toprule\textbf{\normalfont\scriptsize Point de montage / volume} & \textbf{\normalfont\scriptsize Req.} & \textbf{\normalfont\scriptsize Sep.} & \textbf{\normalfont\scriptsize Min} & \textbf{\normalfont\scriptsize Max} & \textbf{\normalfont\scriptsize FS} & \textbf{\normalfont\scriptsize Options de montage} & \textbf{\normalfont\scriptsize Cls} & \textbf{\normalfont\scriptsize Src} & \textbf{\normalfont\scriptsize Notes}\\\midrule\endhead
/ & \textcolor{cBlock}{mandatory} & non & 256 MB & unbounded & JFS2 & rw (default) & system & ibm-release-notes-71 & hd4 logical volume; typically 0.25-2 GB in practice; root filesystem for AIX system\\
\rowcolor{rowalt}
/usr & \textcolor{cBlock}{mandatory} & \textcolor{cOk}{oui} & 1 GB & unbounded & JFS2 & rw (default) & system & ibm-logical-volumes & hd2 logical volume; typically 2.56-6.25 GB; contains operating system binaries and libraries; cannot be moved from rootvg\\
/var & \textcolor{cBlock}{mandatory} & \textcolor{cOk}{oui} & 512 MB & unbounded & JFS2 & rw (default) & system & ibm-logical-volumes & hd9var logical volume; typically 0.56-1 GB; contains variable data like logs and spool; cannot be moved from rootvg\\
\rowcolor{rowalt}
/tmp & \textcolor{cBlock}{mandatory} & \textcolor{cOk}{oui} & 64 MB & unbounded & JFS2 & rw (default) & system & ibm-release-notes-71 & hd3 logical volume; minimum 64 MB enforced during installation/migration; typically 1.4-4 GB; temporary storage; cannot be moved from rootvg\\
/home & \textcolor{cHigh}{recommended} & \textcolor{cOk}{oui} & 1 GB & unbounded & JFS2 & rw (default) & data & ibm-logical-volumes & hd1 logical volume; typically 8-10 GB; user home directories; can be part of rootvg\\
\rowcolor{rowalt}
/opt & \textcolor{cLow}{optional} & conditional & 1 GB & unbounded & JFS2 & rw (default) & system & ibm-logical-volumes & hd10opt logical volume (new in AIX 5.1); optional; only created if no existing /opt filesystem exists; can contain optional software\\
/admin & \textcolor{cLow}{optional} & \textcolor{cOk}{oui} & 512 MB & unbounded & JFS2 & rw (default) & system & ibm-logical-volumes & hd11admin logical volume; optional administrative files and tools filesystem\\
\rowcolor{rowalt}
/var/adm/ras/livedump & \textcolor{cHigh}{recommended} & \textcolor{cOk}{oui} & 64 MB & unbounded & JFS2 & rw (default) & system & ibm-livedumps & livedump logical volume; default size 64 MB or 1/4 system memory (whichever is less); introduced in AIX 6.1; stores livedump crash data\\
N/A (boot) & \textcolor{cBlock}{mandatory} & \textcolor{cOk}{oui} & 24 MB & 24 MB & JFS2 & N/A & boot & ibm-release-notes-71 & hd5 boot logical volume; must be exactly 24 MB for AIX 7.1; must be contiguous and within first 4 GB of disk; bootable kernel and boot files\\
\rowcolor{rowalt}
N/A (paging) & \textcolor{cBlock}{mandatory} & \textcolor{cOk}{oui} & 512 MB & 64 GB (per LV) & paging & N/A & swap & ibm-paging-spaces & hd6 primary paging space logical volume; default 512 MB; recommendation: 2x RAM for RAM $\le$ 256 MB; for RAM > 256 MB: 512 MB + (RAM - 256 MB) x 1.25\\
N/A (journal log) & \textcolor{cBlock}{mandatory} & \textcolor{cOk}{oui} & 4 MB & 256 MB & JFS2 log & N/A & system & ibm-jfs2-log-device & hd8 JFS2 log device; default 1 physical partition (\textasciitilde{}4 MB); IBM recommends 2 MB per 1 GB of data protected; max 256 MB; typically INLINE with filesystem in modern AIX\\
\rowcolor{rowalt}
N/A (dump) & \textcolor{cBlock}{mandatory} & \textcolor{cOk}{oui} & 1.5x RAM & unbounded & dump & N/A & system & ibm-dump-devices & lg\_dumplv (for systems with $\ge$ 4 GB memory) or hd6 (default). lg\_dumplv recommended size: 1.5x estimated system memory for fluctuations. Stores system crash dumps.\\
datavg (application) & \textcolor{cHigh}{recommended} & conditional & unbounded & unbounded & JFS2 & rw (default) & data & ibm-datavg-best-practice & Separate user-defined volume group (datavg) for application data and software; best practice to keep application data out of rootvg; enables easier backup, migration, and system cloning\\
\bottomrule\end{longtable}}
\subsection{7.2}
\noindent\small \textbf{FS par d\'efaut :} \texttt{JFS2} \quad\textbullet\quad \textbf{Swap/paging :} The general recommendation is that the sum of the sizes of the paging spaces should be equal to at least twice the size of real memory, up to 256 MB of memory (512 MB paging). For systems with memory > 256 MB: total paging space = 512 MB + (memory size - 256 MB) * 1.25. Each single paging space cannot exceed 64 GB. Default hd6 is 512 MB.\normalsize\par\smallskip
\noindent\small\textit{AIX 7.2 requires minimum 20 GB disk for default installation (includes all devices, Graphics bundle, System Management Client bundle). All standard filesystems default to JFS2. The rootvg contains mandatory system filesystems (hd2, hd3, hd4, hd5, hd6, hd8, hd9var). Optional filesystems include hd1, hd10opt, hd11admin. All rootvg filesystems except hd10opt cannot be moved. crfs command by default creates filesystems with INLINE log device. Application data should reside in separate datavg volume group. Encryption can be selected for hd1, hd2, hd3, hd9var, hd10opt, hd11admin, and dumplv during installation.}\normalsize\par\smallskip
{\renewcommand{\arraystretch}{1.15}\begin{longtable}{M R S Z Z T O C I N}\toprule\textbf{\normalfont\scriptsize Point de montage / volume} & \textbf{\normalfont\scriptsize Req.} & \textbf{\normalfont\scriptsize Sep.} & \textbf{\normalfont\scriptsize Min} & \textbf{\normalfont\scriptsize Max} & \textbf{\normalfont\scriptsize FS} & \textbf{\normalfont\scriptsize Options de montage} & \textbf{\normalfont\scriptsize Cls} & \textbf{\normalfont\scriptsize Src} & \textbf{\normalfont\scriptsize Notes}\\\midrule\endhead
/ & \textcolor{cBlock}{mandatory} & non & 256 MB & unbounded & JFS2 & rw (default) & system & ibm-release-notes-72 & hd4 logical volume; typically 0.25-18 GB in practice; root filesystem for AIX system\\
\rowcolor{rowalt}
/usr & \textcolor{cBlock}{mandatory} & \textcolor{cOk}{oui} & 1 GB & unbounded & JFS2 & rw (default) & system & ibm-logical-volumes & hd2 logical volume; typically 2.56-6.25 GB; contains operating system binaries and libraries; cannot be moved from rootvg\\
/var & \textcolor{cBlock}{mandatory} & \textcolor{cOk}{oui} & 512 MB & unbounded & JFS2 & rw (default) & system & ibm-logical-volumes & hd9var logical volume; typically 0.56-4 GB; contains variable data like logs and spool; cannot be moved from rootvg; may need 512 MB for prerequisite checks\\
\rowcolor{rowalt}
/tmp & \textcolor{cBlock}{mandatory} & \textcolor{cOk}{oui} & 64 MB & unbounded & JFS2 & rw (default) & system & ibm-release-notes-72 & hd3 logical volume; minimum 64 MB enforced during migration; typically 1.4-4.69 GB; temporary storage; cannot be moved from rootvg\\
/home & \textcolor{cHigh}{recommended} & \textcolor{cOk}{oui} & 1 GB & unbounded & JFS2 & rw (default) & data & ibm-logical-volumes & hd1 logical volume; typically 8-64 GB (10 GB default in VIOS); user home directories; can be part of rootvg\\
\rowcolor{rowalt}
/opt & \textcolor{cLow}{optional} & conditional & 1 GB & unbounded & JFS2 & rw (default) & system & ibm-logical-volumes & hd10opt logical volume; optional; only created if no existing /opt filesystem exists; can contain optional software; can be encrypted during installation\\
/admin & \textcolor{cLow}{optional} & \textcolor{cOk}{oui} & 512 MB & unbounded & JFS2 & rw (default) & system & ibm-logical-volumes & hd11admin logical volume; optional administrative files and tools filesystem; can be encrypted during installation\\
\rowcolor{rowalt}
/var/adm/ras/livedump & \textcolor{cHigh}{recommended} & \textcolor{cOk}{oui} & 64 MB & unbounded & JFS2 & rw (default) & system & ibm-livedumps & livedump logical volume; default size 64 MB or 1/4 system memory (whichever is less); stores livedump crash data; dumpcheck runs daily and monitors 25\% threshold\\
N/A (boot) & \textcolor{cBlock}{mandatory} & \textcolor{cOk}{oui} & 24 MB & 24 MB & JFS2 & N/A & boot & ibm-release-notes-72 & hd5 boot logical volume; must be exactly 24 MB for AIX 7.2; must be contiguous and within first 4 GB of disk; bootable kernel and boot files\\
\rowcolor{rowalt}
N/A (paging) & \textcolor{cBlock}{mandatory} & \textcolor{cOk}{oui} & 512 MB & 64 GB (per LV) & paging & N/A & swap & ibm-release-notes-72 & hd6 primary paging space logical volume; default 512 MB for all new/complete overwrite installations; recommendation: 2x RAM for RAM $\le$ 256 MB; for RAM > 256 MB: 512 MB + (RAM - 256 MB) x 1.25\\
N/A (journal log) & \textcolor{cBlock}{mandatory} & \textcolor{cOk}{oui} & 4 MB & 256 MB & JFS2 log & N/A & system & ibm-jfs2-log-device & hd8 JFS2 log device; default 1 physical partition (\textasciitilde{}4 MB); IBM recommends 2 MB per 1 GB of data protected; max 256 MB; crfs by default creates filesystem with INLINE log device\\
\rowcolor{rowalt}
N/A (dump) & \textcolor{cBlock}{mandatory} & \textcolor{cOk}{oui} & 1.5x RAM & unbounded & dump & N/A & system & ibm-dump-devices & lg\_dumplv (for systems with $\ge$ 4 GB memory) or hd6 (default). lg\_dumplv recommended size: 1.5x estimated system memory for fluctuations. Stores system crash dumps.\\
datavg (application) & \textcolor{cHigh}{recommended} & conditional & unbounded & unbounded & JFS2 & rw (default) & data & ibm-datavg-best-practice & Separate user-defined volume group (datavg) for application data and software; best practice to keep application data out of rootvg; enables easier backup, migration, and system cloning\\
\bottomrule\end{longtable}}
\subsection{7.3}
\noindent\small \textbf{FS par d\'efaut :} \texttt{JFS2} \quad\textbullet\quad \textbf{Swap/paging :} The general recommendation is that the sum of the sizes of the paging spaces should be equal to at least twice the size of real memory, up to 256 MB of memory (512 MB paging). For systems with memory > 256 MB: total paging space = 512 MB + (memory size - 256 MB) * 1.25. Each single paging space cannot exceed 64 GB. Default hd6 is 512 MB.\normalsize\par\smallskip
\noindent\small\textit{AIX 7.3 requires minimum 20 GB disk for default installation (includes all devices, Graphics bundle, System Management bundle). All standard filesystems default to JFS2. Key change: hd5 (boot) minimum increased from 24 MB to 40 MB. The rootvg contains mandatory system filesystems (hd2, hd3, hd4, hd5, hd6, hd8, hd9var). Optional filesystems include hd1, hd10opt, hd11admin. All rootvg filesystems except hd10opt cannot be moved. crfs command by default creates filesystems with INLINE log device. Application data should reside in separate datavg volume group. Encryption can be selected for multiple logical volumes during installation.}\normalsize\par\smallskip
{\renewcommand{\arraystretch}{1.15}\begin{longtable}{M R S Z Z T O C I N}\toprule\textbf{\normalfont\scriptsize Point de montage / volume} & \textbf{\normalfont\scriptsize Req.} & \textbf{\normalfont\scriptsize Sep.} & \textbf{\normalfont\scriptsize Min} & \textbf{\normalfont\scriptsize Max} & \textbf{\normalfont\scriptsize FS} & \textbf{\normalfont\scriptsize Options de montage} & \textbf{\normalfont\scriptsize Cls} & \textbf{\normalfont\scriptsize Src} & \textbf{\normalfont\scriptsize Notes}\\\midrule\endhead
/ & \textcolor{cBlock}{mandatory} & non & 256 MB & unbounded & JFS2 & rw (default) & system & ibm-release-notes-73 & hd4 logical volume; typically 0.25-18 GB in practice; root filesystem for AIX system\\
\rowcolor{rowalt}
/usr & \textcolor{cBlock}{mandatory} & \textcolor{cOk}{oui} & 1 GB & unbounded & JFS2 & rw (default) & system & ibm-logical-volumes & hd2 logical volume; typically 2.56-6.25 GB; contains operating system binaries and libraries; cannot be moved from rootvg\\
/var & \textcolor{cBlock}{mandatory} & \textcolor{cOk}{oui} & 512 MB & unbounded & JFS2 & rw (default) & system & ibm-logical-volumes & hd9var logical volume; typically 0.56-4 GB; contains variable data like logs and spool; cannot be moved from rootvg\\
\rowcolor{rowalt}
/tmp & \textcolor{cBlock}{mandatory} & \textcolor{cOk}{oui} & 64 MB & unbounded & JFS2 & rw (default) & system & ibm-release-notes-73 & hd3 logical volume; minimum 64 MB enforced during migration; typically 1.4-4.69 GB; temporary storage; cannot be moved from rootvg\\
/home & \textcolor{cHigh}{recommended} & \textcolor{cOk}{oui} & 1 GB & unbounded & JFS2 & rw (default) & data & ibm-logical-volumes & hd1 logical volume; typically 8-64 GB; user home directories; can be part of rootvg\\
\rowcolor{rowalt}
/opt & \textcolor{cLow}{optional} & conditional & 1 GB & unbounded & JFS2 & rw (default) & system & ibm-logical-volumes & hd10opt logical volume; optional; only created if no existing /opt filesystem exists; can contain optional software; can be encrypted during installation\\
/admin & \textcolor{cLow}{optional} & \textcolor{cOk}{oui} & 512 MB & unbounded & JFS2 & rw (default) & system & ibm-logical-volumes & hd11admin logical volume; optional administrative files and tools filesystem; can be encrypted during installation\\
\rowcolor{rowalt}
/var/adm/ras/livedump & \textcolor{cHigh}{recommended} & \textcolor{cOk}{oui} & 64 MB & unbounded & JFS2 & rw (default) & system & ibm-livedumps & livedump logical volume; default size 64 MB or 1/4 system memory (whichever is less); stores livedump crash data; dumpcheck runs daily and monitors 25\% threshold\\
N/A (boot) & \textcolor{cBlock}{mandatory} & \textcolor{cOk}{oui} & 40 MB & 40 MB & JFS2 & N/A & boot & ibm-release-notes-73 & hd5 boot logical volume; minimum 40 MB for AIX 7.3 (increased from 24 MB in 7.1/7.2); must be contiguous and within first 4 GB of disk; bootable kernel and boot files\\
\rowcolor{rowalt}
N/A (paging) & \textcolor{cBlock}{mandatory} & \textcolor{cOk}{oui} & 512 MB & 64 GB (per LV) & paging & N/A & swap & ibm-release-notes-73 & hd6 primary paging space logical volume; default 512 MB for all new/complete overwrite installations; recommendation: 2x RAM for RAM $\le$ 256 MB; for RAM > 256 MB: 512 MB + (RAM - 256 MB) x 1.25\\
N/A (journal log) & \textcolor{cBlock}{mandatory} & \textcolor{cOk}{oui} & 4 MB & 256 MB & JFS2 log & N/A & system & ibm-jfs2-log-device & hd8 JFS2 log device; default 1 physical partition (\textasciitilde{}4 MB); IBM recommends 2 MB per 1 GB of data protected; max 256 MB; crfs by default creates filesystem with INLINE log device\\
\rowcolor{rowalt}
N/A (dump) & \textcolor{cBlock}{mandatory} & \textcolor{cOk}{oui} & 1.5x RAM & unbounded & dump & N/A & system & ibm-dump-devices & lg\_dumplv (for systems with $\ge$ 4 GB memory) or hd6 (default). lg\_dumplv recommended size: 1.5x estimated system memory for fluctuations. Stores system crash dumps.\\
datavg (application) & \textcolor{cHigh}{recommended} & conditional & unbounded & unbounded & JFS2 & rw (default) & data & ibm-datavg-best-practice & Separate user-defined volume group (datavg) for application data and software; best practice to keep application data out of rootvg; enables easier backup, migration, and system cloning\\
\bottomrule\end{longtable}}
\clearpage
\section{Windows Server}
\noindent\small\textit{This research covers the four major long-term servicing channel (LTSC) releases of Windows Server: 2016, 2019, 2022, and 2025. All partition sizing and filesystem recommendations derive from official Microsoft Learn documentation, with specific guidance from hardware requirements, BitLocker planning, and UEFI/GPT configuration articles. The CIS Benchmark sources referenced are available through the official CIS Security portal, though detailed benchmark recommendations beyond filesystem requirements would require accessing the full benchmark documents (v4.0.0 for 2022 published May 2025, v1.0.0 for 2025 published March 2025). Storage Spaces and ReFS guidance sourced from official Windows Server storage documentation updated through 2026. Pagefile sizing based on official Microsoft performance guidance with automatic memory dump as the default setting across all four versions. NTFS is the consistent default filesystem for system partitions across all versions; ReFS is recommended as an option for data volumes and Storage Spaces Direct deployments, particularly in 2022 and 2025. All mount\_point references follow Windows convention (drive letters C:, D:, etc., or partition names for system partitions). No mount\_options specified for NTFS/ReFS beyond noting that BitLocker, ACLs, and NTFS permissions are applicable; CIS Benchmarks provide additional hardening options not detailed here."
}\normalsize\par\medskip
\subsection{2016}
\noindent\small \textbf{FS par d\'efaut :} \texttt{NTFS} \quad\textbullet\quad \textbf{Swap/paging :} Pagefile sizing depends on crash dump settings and peak system commit charge. For complete memory dump: minimum 1x RAM plus 257 MB. System-managed pagefiles range from RAM/8 (max 32 GB) minimum to 3x RAM or 4 GB maximum (whichever larger). Modern recommendation: size based on actual system committed memory usage, not arbitrary RAM multiplier. Location: C: drive by default.\normalsize\par\smallskip
\noindent\small\textit{Windows Server 2016 uses UEFI/GPT partitioning scheme by default for new installations. All partitions except data can reside on same physical disk. BitLocker requires OS volume on NTFS, system partition on separate partition (FAT32 for UEFI, NTFS for BIOS, min 350 MB). Storage Spaces and ReFS available but NTFS is default for system partition.}\normalsize\par\smallskip
{\renewcommand{\arraystretch}{1.15}\begin{longtable}{M R S Z Z T O C I N}\toprule\textbf{\normalfont\scriptsize Point de montage / volume} & \textbf{\normalfont\scriptsize Req.} & \textbf{\normalfont\scriptsize Sep.} & \textbf{\normalfont\scriptsize Min} & \textbf{\normalfont\scriptsize Max} & \textbf{\normalfont\scriptsize FS} & \textbf{\normalfont\scriptsize Options de montage} & \textbf{\normalfont\scriptsize Cls} & \textbf{\normalfont\scriptsize Src} & \textbf{\normalfont\scriptsize Notes}\\\midrule\endhead
EFI System Partition (ESP) & \textcolor{cBlock}{mandatory} & \textcolor{cOk}{oui} & 200 MB & unbounded & FAT32 & N/A & boot & ms-uefi-gpt-config & Required for UEFI/GPT-based systems. Microsoft documentation recommends 260 MB or higher for reliability. Must be formatted FAT32 for UEFI platforms. System boot files only, no user data.\\
\rowcolor{rowalt}
Microsoft Reserved Partition (MSR) & \textcolor{cBlock}{mandatory} & \textcolor{cOk}{oui} & 16 MB & 128 MB & N/A & N/A & system & ms-uefi-gpt-config & Required for all GPT disks. Default size is 16 MB (minimum since Windows 10). Does not receive a drive letter. Cannot store user data. Used by Windows for partition management.\\
C: (System/Windows Partition) & \textcolor{cBlock}{mandatory} & \textcolor{cOk}{oui} & 20 GB & unbounded & NTFS & N/A (BitLocker ACL enforcement, NTFS permissions) & system & ms-hardware-reqs & Operating system and Windows files. Microsoft minimum is 20 GB for 64-bit versions. 32 GB is recommended as practical minimum for Server Core with full updates. Requires minimum 16 GB free space after OOBE. Must have sufficient space for paging, hibernation, and crash dumps (additional space required if RAM > 16 GB). NTFS is mandatory.\\
\rowcolor{rowalt}
Windows Recovery Environment (WinRE) Partition & \textcolor{cHigh}{recommended} & \textcolor{cOk}{oui} & 300 MB & 1 GB & NTFS & N/A & boot & ms-winre-partition & Minimum 300 MB required. 500-700 MB typical for winre.wim image depending on language/customizations. For Windows Server 2016: minimum 52 MB free space recommended (250 MB best practice). Separate from OS partition for BitLocker failover support. Type ID: DE94BBA4-06D1-4D40-A16A-BFD50179D6AC.\\
D: and beyond (Data Volumes, optional) & \textcolor{cLow}{optional} & conditional & unbounded & 35 PB (ReFS) / 256 TB (NTFS) & NTFS or ReFS & N/A (NTFS ACLs; ReFS for virtualized workloads) & data & ms-volume-planning & Recommended to place after WinRE partition. NTFS supports up to 256 TB volumes. ReFS supports up to 35 PB but recommended limit is 64 TB for Storage Spaces Direct. For file server workloads with VSS backup, limit to 10 TB for performance.\\
\bottomrule\end{longtable}}
\subsection{2019}
\noindent\small \textbf{FS par d\'efaut :} \texttt{NTFS} \quad\textbullet\quad \textbf{Swap/paging :} Pagefile sizing depends on crash dump settings and peak system commit charge. For complete memory dump: minimum 1x RAM plus 257 MB. System-managed pagefiles range from RAM/8 (max 32 GB) minimum to 3x RAM or 4 GB maximum (whichever larger). Modern recommendation: size based on actual system committed memory usage, not arbitrary RAM multiplier. Location: C: drive by default.\normalsize\par\smallskip
\noindent\small\textit{Windows Server 2019 inherits the same partitioning scheme as 2016 with UEFI/GPT as default. BitLocker fully integrated with NTFS volumes. Storage Spaces and ReFS available with improved block cloning and performance features compared to 2016.}\normalsize\par\smallskip
{\renewcommand{\arraystretch}{1.15}\begin{longtable}{M R S Z Z T O C I N}\toprule\textbf{\normalfont\scriptsize Point de montage / volume} & \textbf{\normalfont\scriptsize Req.} & \textbf{\normalfont\scriptsize Sep.} & \textbf{\normalfont\scriptsize Min} & \textbf{\normalfont\scriptsize Max} & \textbf{\normalfont\scriptsize FS} & \textbf{\normalfont\scriptsize Options de montage} & \textbf{\normalfont\scriptsize Cls} & \textbf{\normalfont\scriptsize Src} & \textbf{\normalfont\scriptsize Notes}\\\midrule\endhead
EFI System Partition (ESP) & \textcolor{cBlock}{mandatory} & \textcolor{cOk}{oui} & 200 MB & unbounded & FAT32 & N/A & boot & ms-uefi-gpt-config & Required for UEFI/GPT-based systems. Microsoft documentation recommends 260 MB or higher for reliability. Must be formatted FAT32 for UEFI platforms. System boot files only, no user data.\\
\rowcolor{rowalt}
Microsoft Reserved Partition (MSR) & \textcolor{cBlock}{mandatory} & \textcolor{cOk}{oui} & 16 MB & 128 MB & N/A & N/A & system & ms-uefi-gpt-config & Required for all GPT disks. Default size is 16 MB. Does not receive a drive letter. Cannot store user data. Used by Windows for partition management.\\
C: (System/Windows Partition) & \textcolor{cBlock}{mandatory} & \textcolor{cOk}{oui} & 20 GB & unbounded & NTFS & N/A (BitLocker ACL enforcement, NTFS permissions) & system & ms-hardware-reqs & Operating system and Windows files. Microsoft minimum is 20 GB for 64-bit versions. 32 GB is recommended as practical minimum for Server Core with full updates. Requires minimum 16 GB free space after OOBE. Must have sufficient space for paging, hibernation, and crash dumps (additional space required if RAM > 16 GB). NTFS is mandatory.\\
\rowcolor{rowalt}
Windows Recovery Environment (WinRE) Partition & \textcolor{cHigh}{recommended} & \textcolor{cOk}{oui} & 300 MB & 1 GB & NTFS & N/A & boot & ms-winre-partition & Minimum 300 MB required. 500-700 MB typical for winre.wim image depending on language/customizations. For Windows Server 2019: minimum 52 MB free space recommended (250 MB best practice, same as 2016). Separate from OS partition for BitLocker failover support. Type ID: DE94BBA4-06D1-4D40-A16A-BFD50179D6AC.\\
D: and beyond (Data Volumes, optional) & \textcolor{cLow}{optional} & conditional & unbounded & 35 PB (ReFS) / 256 TB (NTFS) & NTFS or ReFS & N/A (NTFS ACLs; ReFS for virtualized workloads) & data & ms-volume-planning & Recommended to place after WinRE partition. NTFS supports up to 256 TB volumes. ReFS supports up to 35 PB but recommended limit is 64 TB for Storage Spaces Direct. For file server workloads with VSS backup, limit to 10 TB for performance.\\
\bottomrule\end{longtable}}
\subsection{2022}
\noindent\small \textbf{FS par d\'efaut :} \texttt{NTFS} \quad\textbullet\quad \textbf{Swap/paging :} Pagefile sizing depends on crash dump settings and peak system commit charge. For complete memory dump: minimum 1x RAM plus 257 MB. System-managed pagefiles range from RAM/8 (max 32 GB) minimum to 3x RAM or 4 GB maximum (whichever larger). Modern recommendation: size based on actual system committed memory usage, not arbitrary RAM multiplier. Location: C: drive by default. Automatic memory dump is default setting.\normalsize\par\smallskip
\noindent\small\textit{Windows Server 2022 introduces ReFS as recommended for Storage Spaces Direct (not default for system partition). WinRE free space requirement increased to 100 MB minimum (vs. 52 MB in earlier versions). Data Deduplication available on both NTFS and ReFS. File-level snapshots available with ReFS. CIS Benchmark v4.0.0 (May 2025) available for hardening guidance.}\normalsize\par\smallskip
{\renewcommand{\arraystretch}{1.15}\begin{longtable}{M R S Z Z T O C I N}\toprule\textbf{\normalfont\scriptsize Point de montage / volume} & \textbf{\normalfont\scriptsize Req.} & \textbf{\normalfont\scriptsize Sep.} & \textbf{\normalfont\scriptsize Min} & \textbf{\normalfont\scriptsize Max} & \textbf{\normalfont\scriptsize FS} & \textbf{\normalfont\scriptsize Options de montage} & \textbf{\normalfont\scriptsize Cls} & \textbf{\normalfont\scriptsize Src} & \textbf{\normalfont\scriptsize Notes}\\\midrule\endhead
EFI System Partition (ESP) & \textcolor{cBlock}{mandatory} & \textcolor{cOk}{oui} & 200 MB & unbounded & FAT32 & N/A & boot & ms-uefi-gpt-config & Required for UEFI/GPT-based systems. Microsoft documentation recommends 260 MB or higher for reliability. Must be formatted FAT32 for UEFI platforms. System boot files only, no user data.\\
\rowcolor{rowalt}
Microsoft Reserved Partition (MSR) & \textcolor{cBlock}{mandatory} & \textcolor{cOk}{oui} & 16 MB & 128 MB & N/A & N/A & system & ms-uefi-gpt-config & Required for all GPT disks. Default size is 16 MB. Does not receive a drive letter. Cannot store user data. Used by Windows for partition management.\\
C: (System/Windows Partition) & \textcolor{cBlock}{mandatory} & \textcolor{cOk}{oui} & 20 GB & unbounded & NTFS & N/A (BitLocker ACL enforcement, NTFS permissions) & system & ms-hardware-reqs & Operating system and Windows files. Microsoft minimum is 20 GB for 64-bit versions. 32 GB is recommended as practical minimum for Server Core with full updates. Requires minimum 16 GB free space after OOBE. Must have sufficient space for paging, hibernation, and crash dumps (additional space required if RAM > 16 GB). NTFS is mandatory.\\
\rowcolor{rowalt}
Windows Recovery Environment (WinRE) Partition & \textcolor{cHigh}{recommended} & \textcolor{cOk}{oui} & 300 MB & 1 GB & NTFS & N/A & boot & ms-winre-partition & Minimum 300 MB required. 500-700 MB typical for winre.wim image. For Windows Server 2022: minimum 100 MB free space required (upgraded from 52 MB in earlier versions), 250 MB recommended for future updates. Separate from OS partition for BitLocker failover support. Type ID: DE94BBA4-06D1-4D40-A16A-BFD50179D6AC.\\
D: and beyond (Data Volumes, optional) & \textcolor{cLow}{optional} & conditional & unbounded & 35 PB (ReFS) / 256 TB (NTFS) & NTFS or ReFS & N/A (NTFS ACLs; ReFS for virtualized workloads) & data & ms-volume-planning & Recommended to place after WinRE partition. NTFS supports up to 256 TB volumes. ReFS is now recommended for Storage Spaces Direct deployments (available as option). ReFS supports up to 35 PB but recommended limit is 64 TB for Storage Spaces Direct. For file server workloads with VSS backup, limit to 10 TB for performance.\\
\bottomrule\end{longtable}}
\subsection{2025}
\noindent\small \textbf{FS par d\'efaut :} \texttt{NTFS} \quad\textbullet\quad \textbf{Swap/paging :} Pagefile sizing depends on crash dump settings and peak system commit charge. For complete memory dump: minimum 1x RAM plus 257 MB. System-managed pagefiles range from RAM/8 (max 32 GB) minimum to 3x RAM or 4 GB maximum (whichever larger). Modern recommendation: size based on actual system committed memory usage, not arbitrary RAM multiplier. Location: C: drive by default. Automatic memory dump is default setting.\normalsize\par\smallskip
\noindent\small\textit{Windows Server 2025 released November 2024. Introduces native NVMe support with reduced latency (no SCSI translation required), SMB compression (LZ4), and thin provisioned Storage Spaces volumes. WinRE free space guidance increased to 1 GB recommended (200 MB minimum). ReFS continues as recommended for Storage Spaces Direct with enhanced performance. CIS Benchmark v1.0.0 published March 2025. BitLocker support for ReFS volumes introduced.}\normalsize\par\smallskip
{\renewcommand{\arraystretch}{1.15}\begin{longtable}{M R S Z Z T O C I N}\toprule\textbf{\normalfont\scriptsize Point de montage / volume} & \textbf{\normalfont\scriptsize Req.} & \textbf{\normalfont\scriptsize Sep.} & \textbf{\normalfont\scriptsize Min} & \textbf{\normalfont\scriptsize Max} & \textbf{\normalfont\scriptsize FS} & \textbf{\normalfont\scriptsize Options de montage} & \textbf{\normalfont\scriptsize Cls} & \textbf{\normalfont\scriptsize Src} & \textbf{\normalfont\scriptsize Notes}\\\midrule\endhead
EFI System Partition (ESP) & \textcolor{cBlock}{mandatory} & \textcolor{cOk}{oui} & 200 MB & unbounded & FAT32 & N/A & boot & ms-uefi-gpt-config & Required for UEFI/GPT-based systems. Microsoft documentation recommends 260 MB or higher for reliability. Must be formatted FAT32 for UEFI platforms. System boot files only, no user data.\\
\rowcolor{rowalt}
Microsoft Reserved Partition (MSR) & \textcolor{cBlock}{mandatory} & \textcolor{cOk}{oui} & 16 MB & 128 MB & N/A & N/A & system & ms-uefi-gpt-config & Required for all GPT disks. Default size is 16 MB. Does not receive a drive letter. Cannot store user data. Used by Windows for partition management.\\
C: (System/Windows Partition) & \textcolor{cBlock}{mandatory} & \textcolor{cOk}{oui} & 20 GB & unbounded & NTFS & N/A (BitLocker ACL enforcement, NTFS permissions) & system & ms-hardware-reqs & Operating system and Windows files. Microsoft minimum is 20 GB for 64-bit versions. 32 GB is recommended as practical minimum for Server Core with full updates. Requires minimum 16 GB free space after OOBE. Must have sufficient space for paging, hibernation, and crash dumps (additional space required if RAM > 16 GB). NTFS is mandatory.\\
\rowcolor{rowalt}
Windows Recovery Environment (WinRE) Partition & \textcolor{cHigh}{recommended} & \textcolor{cOk}{oui} & 300 MB & 1 GB & NTFS & N/A & boot & ms-winre-partition & Minimum 300 MB required. 500-700 MB typical for winre.wim image. For Windows Server 2025: minimum 200 MB free space required (updated from earlier versions), 1 GB recommended to handle future updates. Separate from OS partition for BitLocker failover support. Type ID: DE94BBA4-06D1-4D40-A16A-BFD50179D6AC.\\
D: and beyond (Data Volumes, optional) & \textcolor{cLow}{optional} & conditional & unbounded & 35 PB (ReFS) / 256 TB (NTFS) & NTFS or ReFS & N/A (NTFS ACLs; ReFS for virtualized workloads) & data & ms-volume-planning & Recommended to place after WinRE partition. NTFS supports up to 256 TB volumes. ReFS is recommended for Storage Spaces Direct deployments. ReFS supports up to 35 PB but recommended limit is 64 TB for Storage Spaces Direct. For file server workloads with VSS backup, limit to 10 TB for performance. Thin provisioned volumes now supported with Storage Spaces Direct.\\
\bottomrule\end{longtable}}
\clearpage
\section{Cross-distro baseline (FHS/CIS)}
\noindent\small\textit{This reference covers the baseline cross-distro standards defined by FHS 3.0 (the foundational standard for Unix-like filesystem hierarchy), general CIS Benchmark controls applicable across Linux distributions, and vendor-specific guidance from the major enterprise distributions: Red Hat (RHEL 8/9/10, XFS default), Ubuntu/Debian (ext4 default, Debian-specific security practices), and SUSE (Btrfs with subvolumes for root, XFS for additional partitions). Each partition is tied to authoritative sources; sizing is based on official vendor documentation and practical deployment recommendations. Mount options follow CIS hardening standards emphasizing nodev, nosuid, and noexec to limit privilege escalation and code execution attacks. Swap sizing is based on Red Hat's modern RAM-based guidance, which supersedes the historical 2x RAM rule. FHS 3.0 does not prescribe partition sizes, deferring to system administrators-sizes listed represent consensus vendor recommendations. Some Red Hat documentation pages (marked retrieved=false) were behind subscription walls but are cited from publicly visible documentation links; exact values in notes are inferred from published best-practice guides. All URLs are live as of June 2026.}\normalsize\par\medskip
\subsection{FHS 3.0 (Filesystem Hierarchy Standard)}
\noindent\small \textbf{FS par d\'efaut :} \texttt{ext4/XFS (distribution-dependent)} \quad\textbullet\quad \textbf{Swap/paging :} FHS 3.0 does not specify swap sizing. Debian recommends swap equal to RAM (minimum 512 MB). Red Hat's modern guidance (applicable across distributions) for systems with 16-64 GB RAM: 4-8 GB swap. For hibernation, allocate swap >= RAM.\normalsize\par\smallskip
\noindent\small\textit{FHS 3.0 is the reference standard defining the hierarchy and purposes of directories across Unix-like systems. The standard itself does NOT specify minimum/maximum partition sizes-those decisions are left to system administrators and distributions. FHS was republished by FreeDesktop in November 2025. Key principle: /usr and /var are designed to be mounted on separate filesystems; /var is required separate from /usr to allow /usr to be mounted read-only.}\normalsize\par\smallskip
{\renewcommand{\arraystretch}{1.15}\begin{longtable}{M R S Z Z T O C I N}\toprule\textbf{\normalfont\scriptsize Point de montage / volume} & \textbf{\normalfont\scriptsize Req.} & \textbf{\normalfont\scriptsize Sep.} & \textbf{\normalfont\scriptsize Min} & \textbf{\normalfont\scriptsize Max} & \textbf{\normalfont\scriptsize FS} & \textbf{\normalfont\scriptsize Options de montage} & \textbf{\normalfont\scriptsize Cls} & \textbf{\normalfont\scriptsize Src} & \textbf{\normalfont\scriptsize Notes}\\\midrule\endhead
/ & \textcolor{cBlock}{mandatory} & non & unbounded & unbounded & ext4 or XFS & defaults & system & FHS3.0-ch03 & Root filesystem must contain essential files to boot, restore, recover, and repair the system. Contains /etc, /bin, /sbin, /lib, and other critical system files.\\
\rowcolor{rowalt}
/usr & \textcolor{cBlock}{mandatory} & conditional & unbounded & unbounded & ext4 or XFS & defaults,ro (read-only recommended) & system & FHS3.0-ch04 & Secondary hierarchy containing user programs, libraries, and documentation. FHS states /usr should be shareable, read-only data. Optional separate partition recommended for easier management and to mount read-only.\\
/var & \textcolor{cHigh}{recommended} & \textcolor{cOk}{oui} & 3 GiB & unbounded & ext4 or XFS & defaults & data & FHS3.0-ch05,Debian-secure-manual & Variable state information: logs, caches, mail spools, package caches (apt/yum). Separate partition prevents logs from filling root filesystem. Debian recommends larger allocation due to package cache in /var/cache/apt/archives.\\
\rowcolor{rowalt}
/var/log & \textcolor{cHigh}{recommended} & \textcolor{cOk}{oui} & 5 GiB & unbounded & ext4 or XFS & defaults,nodev,nosuid,noexec & data & CIS-Linux-partition,FHS3.0-ch05 & Separate partition protects audit logs and system logs from root filesystem exhaustion. Default systemd journal size limited to 10\% of partition size, capped at 4 GiB. CIS Benchmark mandatory control 1.1.5/1.1.6.\\
/var/log/audit & \textcolor{cHigh}{recommended} & \textcolor{cOk}{oui} & 2 GiB & unbounded & ext4 or XFS & defaults,nodev,nosuid,noexec & data & CIS-Linux-partition,TechGirlKB & Separate partition for audit logs (auditd) protects against resource exhaustion and data loss. CIS Benchmark mandatory control 1.1.6/1.1.7. Mount options prevent device files, SUID bits, and executable binaries.\\
\rowcolor{rowalt}
/var/tmp & \textcolor{cHigh}{recommended} & conditional & 1 GiB & unbounded & ext4 or XFS & defaults,nodev,nosuid,noexec,bind (or symlink to /tmp) & data & CIS-Linux-partition,Debian-secure-manual & Temporary files that persist between reboots (contrast to /tmp). CIS Benchmark requires separate partition OR bind mount from /tmp with nodev,nosuid,noexec. Best practice: bind /var/tmp to /tmp to avoid duplicate filesystem.\\
/tmp & \textcolor{cHigh}{recommended} & \textcolor{cOk}{oui} & 1 GiB & unbounded & ext4 or XFS & defaults,nodev,nosuid,noexec & data & FHS3.0-ch03,CIS-Linux-partition & Temporary file storage cleared between system boots (recommended but not required by FHS). Separate partition prevents application temp files from filling root. CIS Benchmark mandatory control 1.1.2/1.1.3. Mount options prevent SUID binaries and device files.\\
\rowcolor{rowalt}
/home & \textcolor{cHigh}{recommended} & \textcolor{cOk}{oui} & 10 GiB & unbounded & ext4 or XFS & defaults,nodev,nosuid & data & FHS3.0-ch03,CIS-Linux-partition,Debian-secure-manual & User home directories (optional per FHS but strongly recommended). Separate partition allows system reinstallation without data loss. CIS Benchmark requires nodev,nosuid. Debian: large partition recommended for multi-user systems.\\
/boot & \textcolor{cHigh}{recommended} & conditional & 250 MiB & 1 GiB & ext4 or XFS & defaults,nodev,nosuid,noexec & boot & RHEL-install-guide,FHS3.0-ch03,CIS-Linux-partition & Static boot files: kernel images, bootloader data, initramfs. Red Hat: 250 MB sufficient for most systems (each kernel \textasciitilde{}30 MB). UEFI systems may allocate 500MB-1GB. Separate /boot/efi 50-150 MB for UEFI. CIS Benchmark recommends nodev,nosuid,noexec.\\
\rowcolor{rowalt}
/boot/efi & \textcolor{cMed}{conditional} & \textcolor{cOk}{oui} & 50 MiB & 150 MiB & EFI System Partition (FAT32/vfat) & defaults,nodev,nosuid,noexec & boot & RHEL-install-guide & UEFI systems only. EFI firmware partition. Not applicable to legacy BIOS systems.\\
/run & \textcolor{cBlock}{mandatory} & non & unbounded & unbounded & tmpfs & defaults,nosuid,nodev,mode=755 & system & FHS3.0-ch03 & Runtime data describing the system since boot. Must be tmpfs or ramfs. Files cleared at boot. FHS 3.0 defines purpose and clearing requirements.\\
\rowcolor{rowalt}
/dev/shm & \textcolor{cBlock}{mandatory} & non & unbounded & unbounded & tmpfs & defaults,nosuid,nodev,noexec & system & CIS-Linux-partition & Shared memory filesystem. CIS Benchmark mandatory control 1.1.15/1.1.16 requires noexec,nodev,nosuid mount options to prevent device files and executable code in shared memory.\\
/opt & \textcolor{cLow}{optional} & conditional & unbounded & unbounded & ext4 or XFS & defaults & data & FHS3.0-ch06 & Add-on application software packages. Each package uses /opt/<package> or /opt/<provider>. Optional separate partition useful for preserving third-party software during system reinstalls (SUSE recommendation).\\
\rowcolor{rowalt}
/srv & \textcolor{cLow}{optional} & conditional & unbounded & unbounded & ext4 or XFS & defaults & data & FHS3.0-ch07 & Site-specific data served by services (web, FTP, mail). Optional separate partition for service data. SUSE recommends separate mount point to preserve service data during system reinstalls.\\
/usr/local & \textcolor{cLow}{optional} & conditional & unbounded & unbounded & ext4 or XFS & defaults & data & FHS3.0-ch04,Debian-secure-manual & Locally compiled software and manual installations. Optional separate partition if frequent non-package software installations expected. Debian/SUSE recommend separate partition to preserve during system upgrades.\\
\rowcolor{rowalt}
/etc & \textcolor{cBlock}{mandatory} & non & unbounded & unbounded & ext4 or XFS & defaults & system & FHS3.0-ch03 & Host-specific system configuration files (static, non-executable). Must remain on root filesystem. FHS prohibits binaries in /etc; scripts permitted.\\
\bottomrule\end{longtable}}
\subsection{CIS Benchmark for Linux (General Framework)}
\noindent\small \textbf{FS par d\'efaut :} \texttt{ext4 or XFS (distribution-dependent)} \quad\textbullet\quad \textbf{Swap/paging :} CIS does not explicitly specify swap sizing. Refer to distribution guidance: Red Hat recommends 4-8 GB for systems with 16-64 GB RAM. For hibernation, swap >= RAM. Swap should be encrypted on systems with sensitive data.\normalsize\par\smallskip
\noindent\small\textit{CIS Benchmarks (Center for Internet Security) define security hardening controls for Linux distributions. Section 1.1 covers filesystem configuration with mandatory separate partitions and mount options. Applies to RHEL, Ubuntu, Debian, CentOS, Oracle Linux, etc. CIS controls are risk-based security requirements, not default installations.}\normalsize\par\smallskip
{\renewcommand{\arraystretch}{1.15}\begin{longtable}{M R S Z Z T O C I N}\toprule\textbf{\normalfont\scriptsize Point de montage / volume} & \textbf{\normalfont\scriptsize Req.} & \textbf{\normalfont\scriptsize Sep.} & \textbf{\normalfont\scriptsize Min} & \textbf{\normalfont\scriptsize Max} & \textbf{\normalfont\scriptsize FS} & \textbf{\normalfont\scriptsize Options de montage} & \textbf{\normalfont\scriptsize Cls} & \textbf{\normalfont\scriptsize Src} & \textbf{\normalfont\scriptsize Notes}\\\midrule\endhead
/tmp & \textcolor{cBlock}{mandatory} & \textcolor{cOk}{oui} & 1 GiB & unbounded & ext4 or XFS & defaults,nodev,nosuid,noexec & data & CIS-Linux-1.1.2 & CIS Control 1.1.2/1.1.3: Mandatory separate partition. Mount options prevent device files, SUID binaries, and executable code. nosuid prevents privilege escalation via SUID; noexec prevents binary execution; nodev prevents device file access.\\
\rowcolor{rowalt}
/var & \textcolor{cBlock}{mandatory} & \textcolor{cOk}{oui} & 5 GiB & unbounded & ext4 or XFS & defaults,nodev,nosuid & data & CIS-Linux-1.1.4 & CIS Control 1.1.4/1.1.5: Mandatory separate partition. nodev prevents device files; nosuid prevents SUID bits. Protects logs and cache from filling root filesystem.\\
/var/log & \textcolor{cBlock}{mandatory} & \textcolor{cOk}{oui} & 5 GiB & unbounded & ext4 or XFS & defaults,nodev,nosuid,noexec & data & CIS-Linux-1.1.5 & CIS Control 1.1.5/1.1.6: Mandatory separate partition. Mount options prevent device access and code execution. Protects logs from being overwritten or modified via /dev/* or binary execution.\\
\rowcolor{rowalt}
/var/log/audit & \textcolor{cBlock}{mandatory} & \textcolor{cOk}{oui} & 2 GiB & unbounded & ext4 or XFS & defaults,nodev,nosuid,noexec & data & CIS-Linux-1.1.6 & CIS Control 1.1.6/1.1.7: Mandatory separate partition. auditd logs must be protected from root-level disk-space exhaustion and unauthorized modification. Audit data is critical for compliance and incident investigation.\\
/var/tmp & \textcolor{cBlock}{mandatory} & conditional & 1 GiB & unbounded & ext4 or XFS & defaults,nodev,nosuid,noexec & data & CIS-Linux-1.1.7 & CIS Control 1.1.7/1.1.8: Separate partition OR bind mount from /tmp. Mount options prevent device files, SUID, and executables. /var/tmp persists across reboots (unlike /tmp); separate partition or bind mount prevents temp-file attacks.\\
\rowcolor{rowalt}
/home & \textcolor{cBlock}{mandatory} & \textcolor{cOk}{oui} & 10 GiB & unbounded & ext4 or XFS & defaults,nodev,nosuid & data & CIS-Linux-1.1.9 & CIS Control 1.1.9/1.1.10: Mandatory separate partition. nosuid prevents user-created SUID binaries; nodev prevents device file attacks. Isolates user data and prevents denial-of-service via home directory expansion.\\
/dev/shm & \textcolor{cBlock}{mandatory} & non & unbounded & unbounded & tmpfs & defaults,nodev,nosuid,noexec & system & CIS-Linux-1.1.15 & CIS Control 1.1.15/1.1.16/1.1.17: tmpfs shared memory mount. noexec prevents executable code; nodev blocks device files; nosuid blocks SUID. Protects IPC shared memory from code injection.\\
\rowcolor{rowalt}
/boot & \textcolor{cHigh}{recommended} & conditional & 250 MiB & 1 GiB & ext4 or XFS & defaults,nodev,nosuid,noexec & boot & CIS-Linux-hardening & CIS general recommendation (not all versions mandate). Mount options: nodev prevents device creation; nosuid prevents SUID; noexec prevents binary execution. Protects kernel and bootloader from modification.\\
\bottomrule\end{longtable}}
\subsection{Red Hat Enterprise Linux (RHEL 8, 9, 10) - Vendor Guidance}
\noindent\small \textbf{FS par d\'efaut :} \texttt{XFS} \quad\textbullet\quad \textbf{Swap/paging :} Red Hat guidance (RHEL 8/9/10): RAM <= 4 GB: minimum 2 GB swap. 4-16 GB RAM: 4 GB swap. 16-64 GB RAM: 8 GB swap. 64-256 GB RAM: 16 GB swap. For hibernation: swap >= RAM. Modern systems with high RAM rarely need swap equal to RAM.\normalsize\par\smallskip
\noindent\small\textit{RHEL 8, 9, 10 default to XFS filesystem for all partitions except /boot/efi (vfat). XFS is recommended for production systems, large files, and high I/O concurrency. ext4 supported as alternative for smaller systems. RHEL provides detailed installation guidance; CIS Benchmark controls apply for hardening.}\normalsize\par\smallskip
{\renewcommand{\arraystretch}{1.15}\begin{longtable}{M R S Z Z T O C I N}\toprule\textbf{\normalfont\scriptsize Point de montage / volume} & \textbf{\normalfont\scriptsize Req.} & \textbf{\normalfont\scriptsize Sep.} & \textbf{\normalfont\scriptsize Min} & \textbf{\normalfont\scriptsize Max} & \textbf{\normalfont\scriptsize FS} & \textbf{\normalfont\scriptsize Options de montage} & \textbf{\normalfont\scriptsize Cls} & \textbf{\normalfont\scriptsize Src} & \textbf{\normalfont\scriptsize Notes}\\\midrule\endhead
/ & \textcolor{cBlock}{mandatory} & non & 15 GiB & unbounded & XFS (or ext4) & defaults,relatime & system & RHEL-install-guide & RHEL default root partition size 15-20 GiB sufficient for most installations; desktop systems 25-50 GiB. XFS is default and recommended. relatime balances access-time updates with performance.\\
\rowcolor{rowalt}
/boot & \textcolor{cBlock}{mandatory} & \textcolor{cOk}{oui} & 250 MiB & 1 GiB & XFS (or ext4) & defaults,nodev,nosuid,noexec,relatime & boot & RHEL-install-guide & RHEL: 250 MB sufficient for most; each kernel \textasciitilde{}30 MB. UEFI systems require separate /boot/efi 50-150 MB vfat partition. Standard partition, not LVM.\\
/var & \textcolor{cHigh}{recommended} & \textcolor{cOk}{oui} & 5 GiB & unbounded & XFS (or ext4) & defaults,relatime & data & RHEL-install-guide & RHEL: 3-5 GiB minimum; production systems often 20-30 GiB. Contains logs, cache, spools. Separate partition recommended to prevent root exhaustion.\\
\rowcolor{rowalt}
/tmp & \textcolor{cHigh}{recommended} & \textcolor{cOk}{oui} & 1 GiB & unbounded & XFS (or ext4) & defaults,nodev,nosuid,noexec,relatime & data & RHEL-install-guide & RHEL: 1-5 GiB typical; production systems may need larger. Separate partition prevents application temp-file DoS.\\
/home & \textcolor{cHigh}{recommended} & \textcolor{cOk}{oui} & 5 GiB & unbounded & XFS (or ext4) & defaults,nodev,nosuid,relatime & data & RHEL-install-guide & RHEL: 5-10 GiB minimum; production systems 20+ GiB. Separate partition allows system reinstall without user data loss.\\
\bottomrule\end{longtable}}
\subsection{Ubuntu \& Debian - Vendor Guidance}
\noindent\small \textbf{FS par d\'efaut :} \texttt{ext4} \quad\textbullet\quad \textbf{Swap/paging :} Debian: swap equal to RAM (minimum 512 MB). Ubuntu: modern systems use 2 GB swap file by default (not partition). For hibernation: swap >= RAM + headroom for in-use swap. Cloud images may omit swap entirely.\normalsize\par\smallskip
\noindent\small\textit{Ubuntu and Debian default to ext4 filesystem. Ubuntu installer offers options for LVM and encryption. Debian Installer follows FHS and Securing Debian Manual recommendations. CIS Benchmarks for Ubuntu 20.04/22.04/24.04 align with partition controls but Debian/Ubuntu favor ext4 over XFS for desktop/server balance.}\normalsize\par\smallskip
{\renewcommand{\arraystretch}{1.15}\begin{longtable}{M R S Z Z T O C I N}\toprule\textbf{\normalfont\scriptsize Point de montage / volume} & \textbf{\normalfont\scriptsize Req.} & \textbf{\normalfont\scriptsize Sep.} & \textbf{\normalfont\scriptsize Min} & \textbf{\normalfont\scriptsize Max} & \textbf{\normalfont\scriptsize FS} & \textbf{\normalfont\scriptsize Options de montage} & \textbf{\normalfont\scriptsize Cls} & \textbf{\normalfont\scriptsize Src} & \textbf{\normalfont\scriptsize Notes}\\\midrule\endhead
/ & \textcolor{cBlock}{mandatory} & non & 8 GiB & unbounded & ext4 & defaults,relatime,errors=remount-ro & system & Debian-install-guide,Ubuntu-partitioning & Debian: 8 GiB minimum. Ubuntu: 15-25 GiB recommended for desktop, 50+ GiB for servers. errors=remount-ro standard Debian protection.\\
\rowcolor{rowalt}
/boot & \textcolor{cLow}{optional} & conditional & 250 MiB & 500 MiB & ext4 & defaults,nodev,nosuid,noexec,relatime & boot & Ubuntu-partitioning,Debian-install-guide & Ubuntu: "DO NOT create separate /boot" for typical home installations (unnecessary, can cause package-manager issues). Debian: optional. If separate, 250-500 MB sufficient.\\
/var & \textcolor{cHigh}{recommended} & \textcolor{cOk}{oui} & 3 GiB & unbounded & ext4 & defaults,relatime & data & Debian-install-guide,Debian-secure-manual & Debian: 3+ GiB for multi-user systems. Debian larger due to apt cache (/var/cache/apt/archives). Separate partition prevents root exhaustion.\\
\rowcolor{rowalt}
/tmp & \textcolor{cHigh}{recommended} & \textcolor{cOk}{oui} & 1 GiB & unbounded & ext4 (or ext2 for Debian without journaling) & defaults,nodev,nosuid,noexec,relatime & data & Debian-install-guide,CIS-Ubuntu & Debian: optional but recommended. Debian notes ext2 acceptable for /tmp (cleaned at boot). Ubuntu: 5+ GiB practical. CIS requires nodev,nosuid,noexec.\\
/home & \textcolor{cHigh}{recommended} & \textcolor{cOk}{oui} & 10 GiB & unbounded & ext4 & defaults,nodev,nosuid,relatime & data & Ubuntu-partitioning,Debian-secure-manual,CIS-Ubuntu & Ubuntu: highly recommended to preserve personal files during reinstall. Desktop: 30-100+ GiB. Multi-user systems: proportional to user count. CIS: nosuid required.\\
\bottomrule\end{longtable}}
\subsection{SUSE Linux Enterprise Server (SLES 15.x) - Vendor Guidance}
\noindent\small \textbf{FS par d\'efaut :} \texttt{Btrfs (root), XFS (additional partitions)} \quad\textbullet\quad \textbf{Swap/paging :} SUSE: swap >= 256 MB minimum. For hibernation: 512 MB - 1 GB. Modern systems with abundant RAM need modest swap; historical 2x RAM rule outdated.\normalsize\par\smallskip
\noindent\small\textit{SUSE defaults to Btrfs with snapshots for root partition (modern rollback capability) and XFS for additional partitions. Btrfs subvolumes (/, /home, /opt, /srv, /var, /usr/local) provide snapshot isolation without separate filesystems. SUSE Expert Partitioner offers flexible configuration.}\normalsize\par\smallskip
{\renewcommand{\arraystretch}{1.15}\begin{longtable}{M R S Z Z T O C I N}\toprule\textbf{\normalfont\scriptsize Point de montage / volume} & \textbf{\normalfont\scriptsize Req.} & \textbf{\normalfont\scriptsize Sep.} & \textbf{\normalfont\scriptsize Min} & \textbf{\normalfont\scriptsize Max} & \textbf{\normalfont\scriptsize FS} & \textbf{\normalfont\scriptsize Options de montage} & \textbf{\normalfont\scriptsize Cls} & \textbf{\normalfont\scriptsize Src} & \textbf{\normalfont\scriptsize Notes}\\\midrule\endhead
/ & \textcolor{cBlock}{mandatory} & non & 16 GiB & unbounded & Btrfs (with snapshots) & defaults,relatime,subvol=@ & system & SUSE-SLES-expert-partitioner & SUSE default Btrfs with snapshots: minimum 16 GiB, recommended 32+ GiB. Btrfs subvolumes provide copy-on-write and snapshot isolation. Supports online growing, multi-device spanning.\\
\rowcolor{rowalt}
/home & \textcolor{cHigh}{recommended} & \textcolor{cOk}{oui} & 10 GiB & unbounded & Btrfs subvolume or XFS & defaults,relatime,subvol=@home & data & SUSE-SLES-expert-partitioner & SUSE recommendation: separate Btrfs subvolume to protect user data during system rollback. Prevents data loss during OS updates/reinstalls.\\
/var & \textcolor{cHigh}{recommended} & \textcolor{cOk}{oui} & 5 GiB & unbounded & Btrfs subvolume or XFS & defaults,relatime,subvol=@var & data & SUSE-SLES-expert-partitioner & SUSE recommendation: separate Btrfs subvolume to exclude logs from snapshots and preserve across rollbacks. Prevents snapshot bloat from large log files.\\
\rowcolor{rowalt}
/opt & \textcolor{cLow}{optional} & conditional & unbounded & unbounded & Btrfs subvolume or XFS & defaults,relatime,subvol=@opt & data & SUSE-SLES-expert-partitioner & SUSE recommendation: separate Btrfs subvolume to preserve third-party applications during system rollback.\\
/srv & \textcolor{cLow}{optional} & conditional & unbounded & unbounded & Btrfs subvolume or XFS & defaults,relatime,subvol=@srv & data & SUSE-SLES-expert-partitioner & SUSE recommendation: separate Btrfs subvolume to protect service data during system updates.\\
\rowcolor{rowalt}
/usr/local & \textcolor{cLow}{optional} & conditional & unbounded & unbounded & Btrfs subvolume or XFS & defaults,relatime,subvol=@usr\_local & data & SUSE-SLES-expert-partitioner & SUSE recommendation: separate Btrfs subvolume to preserve manually installed software during system rollback.\\
\bottomrule\end{longtable}}
\clearpage
\section{Sources}
\subsection{RHEL family}
\begin{itemize}[leftmargin=1.6em]
\item \textbf{[rhel7-install-guide]} Installation Destination | Installation Guide | Red Hat Enterprise Linux 7 --- \textit{Red Hat} (RHEL 7 (2013-2024)) \textcolor{cOk}{[consult\'e]}\\ \url{https://docs.redhat.com/en/documentation/red_hat_enterprise_linux/7/html/installation_guide/sect-disk-partitioning-setup-x86}
\item \textbf{[rhel7-security-guide]} Partitioning the Disk | Security Guide | Red Hat Enterprise Linux 7 --- \textit{Red Hat} (RHEL 7 (2013-2024)) \textcolor{cMed}{[non v\'erifi\'e]}\\ \url{https://docs.redhat.com/en/documentation/red_hat_enterprise_linux/7/html/security_guide/sec-partitioning_the_disk}
\item \textbf{[rhel7-storage-admin]} Swap Space | Storage Administration Guide | Red Hat Enterprise Linux 7 --- \textit{Red Hat} (RHEL 7 (2013-2024)) \textcolor{cMed}{[non v\'erifi\'e]}\\ \url{https://docs.redhat.com/en/documentation/red_hat_enterprise_linux/7/html/storage_administration_guide/ch-swapspace}
\item \textbf{[rhel8-install-guide]} Partitioning Reference | Performing a Standard RHEL 8 Installation --- \textit{Red Hat} (RHEL 8 (2019-2024)) \textcolor{cMed}{[non v\'erifi\'e]}\\ \url{https://docs.redhat.com/en/documentation/red_hat_enterprise_linux/8/html/performing_a_standard_rhel_8_installation/partitioning-reference_installing-rhel}
\item \textbf{[rhel8-security-hardening]} Security Hardening | Red Hat Enterprise Linux 8 --- \textit{Red Hat} (RHEL 8 (2019-2024)) \textcolor{cMed}{[non v\'erifi\'e]}\\ \url{https://docs.redhat.com/en/documentation/red_hat_enterprise_linux/8/html/security_hardening/assembly_securing-rhel-during-installation-security-hardening}
\item \textbf{[rhel8-managing-storage]} Getting Started with Swap | Managing Storage Devices | Red Hat Enterprise Linux 8 --- \textit{Red Hat} (RHEL 8 (2019-2024)) \textcolor{cMed}{[non v\'erifi\'e]}\\ \url{https://docs.redhat.com/en/documentation/red_hat_enterprise_linux/8/html/managing_storage_devices/getting-started-with-swap_managing-storage-devices}
\item \textbf{[rhel9-install-guide]} Partitioning Reference | Performing a Standard RHEL 9 Installation --- \textit{Red Hat} (RHEL 9 (2022-2024)) \textcolor{cMed}{[non v\'erifi\'e]}\\ \url{https://docs.redhat.com/en/documentation/red_hat_enterprise_linux/9/html/performing_a_standard_rhel_9_installation/partitioning-reference_installing-rhel}
\item \textbf{[rhel9-security-hardening]} Security Hardening | Red Hat Enterprise Linux 9 --- \textit{Red Hat} (RHEL 9 (2022-2024)) \textcolor{cMed}{[non v\'erifi\'e]}\\ \url{https://docs.redhat.com/en/documentation/red_hat_enterprise_linux/9/html/security_hardening/assembly_securing-rhel-during-installation-security-hardening}
\item \textbf{[rhel9-managing-storage]} Getting Started with Swap | Managing Storage Devices | Red Hat Enterprise Linux 9 --- \textit{Red Hat} (RHEL 9 (2022-2024)) \textcolor{cMed}{[non v\'erifi\'e]}\\ \url{https://docs.redhat.com/en/documentation/red_hat_enterprise_linux/9/html/managing_storage_devices/getting-started-with-swap_managing-storage-devices}
\item \textbf{[cis-rhel7]} CIS Red Hat Enterprise Linux 7 Benchmark --- \textit{Center for Internet Security (CIS)} (v4.0.0 (2023-12-21)) \textcolor{cMed}{[non v\'erifi\'e]}\\ \url{https://www.cisecurity.org/benchmark/red_hat_linux}
\item \textbf{[cis-rhel8]} CIS Red Hat Enterprise Linux 8 Benchmark --- \textit{Center for Internet Security (CIS)} (v2.0.0) \textcolor{cMed}{[non v\'erifi\'e]}\\ \url{https://www.cisecurity.org/benchmark/red_hat_linux}
\item \textbf{[cis-rhel9]} CIS Red Hat Enterprise Linux 9 Benchmark --- \textit{Center for Internet Security (CIS)} (v2.0.0 (2024-06-24)) \textcolor{cMed}{[non v\'erifi\'e]}\\ \url{https://www.cisecurity.org/benchmark/red_hat_linux}
\item \textbf{[fhs-3.0]} Filesystem Hierarchy Standard (FHS) 3.0 --- \textit{The Linux Foundation} (3.0 (2015-06-03)) \textcolor{cMed}{[non v\'erifi\'e]}\\ \url{https://refspecs.linuxfoundation.org/FHS_3.0/fhs-3.0.pdf}
\item \textbf{[rhel7-73-release-notes]} Installation and Booting | Red Hat Enterprise Linux 7.3 Release Notes --- \textit{Red Hat} (RHEL 7.3 (2016)) \textcolor{cOk}{[consult\'e]}\\ \url{https://docs.redhat.com/en/documentation/red_hat_enterprise_linux/7/html/7.3_release_notes/bug_fixes_installation_and_booting}
\item \textbf{[rhel-swap-sizing]} Recommended swap size for Red Hat platforms --- \textit{Red Hat} (2024) \textcolor{cMed}{[non v\'erifi\'e]}\\ \url{https://access.redhat.com/solutions/15244}
\item \textbf{[mount-options-effects]} Effects of changing mount options with nodev,noexec,nosuid on /tmp and /dev/shm --- \textit{Red Hat} (2024) \textcolor{cMed}{[non v\'erifi\'e]}\\ \url{https://access.redhat.com/solutions/223643}
\item \textbf{[rhel-var-log-audit-separate]} How to move /var/log/audit to a separate filesystem --- \textit{Red Hat} (2024) \textcolor{cMed}{[non v\'erifi\'e]}\\ \url{https://access.redhat.com/articles/7025414}
\end{itemize}
\subsection{Debian/Ubuntu}
\begin{itemize}[leftmargin=1.6em]
\item \textbf{[debian-install-guide]} Debian GNU/Linux Installation Guide - Appendix C: Partitioning for Debian --- \textit{Debian Project} (Current (Bookworm 12)) \textcolor{cOk}{[consult\'e]}\\ \url{https://www.debian.org/releases/stable/amd64/apc.en.html}
\item \textbf{[ubuntu-disk-space]} Ubuntu Community Help: DiskSpace --- \textit{Ubuntu Community} (Current) \textcolor{cOk}{[consult\'e]}\\ \url{https://help.ubuntu.com/community/DiskSpace}
\item \textbf{[ubuntu-partitioning]} Ubuntu Community Help: PartitioningSchemes --- \textit{Ubuntu Community} (Current) \textcolor{cOk}{[consult\'e]}\\ \url{https://help.ubuntu.com/community/PartitioningSchemes}
\item \textbf{[ubuntu-swap-faq]} Ubuntu Community Help: SwapFaq --- \textit{Ubuntu Community} (Current) \textcolor{cOk}{[consult\'e]}\\ \url{https://help.ubuntu.com/community/SwapFaq}
\item \textbf{[cis-debian-11]} CIS Debian Linux 11 Benchmark v1.0.0 --- \textit{Center for Internet Security (CIS)} (v1.0.0 (2022-09-22)) \textcolor{cMed}{[non v\'erifi\'e]}\\ \url{https://www.itsecure.hu/wp-content/uploads/2023/12/CIS_Debian_Linux_11_Benchmark_v1.0.0.pdf}
\item \textbf{[cis-debian-12]} CIS Debian Linux 12 Benchmark v1.1.0 --- \textit{Center for Internet Security (CIS)} (v1.1.0 (2024-09-26)) \textcolor{cMed}{[non v\'erifi\'e]}\\ \url{https://itsecure.hu/wp-content/uploads/2025/05/CIS_Debian_Linux_12_Benchmark_v1.1.0.pdf}
\item \textbf{[cis-ubuntu-20]} CIS Ubuntu Linux 20.04 LTS Benchmark v1.1.0 --- \textit{Center for Internet Security (CIS)} (v1.1.0 (2020)) \textcolor{cMed}{[non v\'erifi\'e]}\\ \url{https://edu.anarcho-copy.org/GNU%20Linux%20-%20Unix-Like/Debian/CIS_Ubuntu_Linux_20.04_LTS_Benchmark_v1.1.0.pdf}
\item \textbf{[cis-ubuntu-22]} CIS Ubuntu Linux 22.04 LTS Benchmark v1.0.0 --- \textit{Center for Internet Security (CIS)} (v1.0.0 (2022-08-30)) \textcolor{cMed}{[non v\'erifi\'e]}\\ \url{https://www.cisecurity.org/benchmark/ubuntu_linux}
\item \textbf{[cis-ubuntu-24]} CIS Ubuntu Linux 24.04 LTS Benchmark v1.0.0 --- \textit{Center for Internet Security (CIS)} (v1.0.0 (2024)) \textcolor{cMed}{[non v\'erifi\'e]}\\ \url{https://www.scribd.com/document/798893779/CIS-Ubuntu-Linux-24-04-LTS-Benchmark-v1-0-0}
\item \textbf{[fhs-3.0]} Filesystem Hierarchy Standard Version 3.0 --- \textit{Linux Foundation (LSB Workgroup)} (3.0 (2015-06-17)) \textcolor{cOk}{[consult\'e]}\\ \url{https://refspecs.linuxfoundation.org/FHS_3.0/fhs-3.0.html}
\item \textbf{[debian-securing-manual]} Securing Debian Manual - Chapter 4.10: Mounting partitions the right way --- \textit{Debian Project} (Current) \textcolor{cOk}{[consult\'e]}\\ \url{https://www.debian.org/doc/manuals/securing-debian-manual/ch04s10.en.html}
\item \textbf{[debian-wiki-partition]} Debian Wiki: Partition --- \textit{Debian Community} (Current) \textcolor{cOk}{[consult\'e]}\\ \url{https://wiki.debian.org/Partition}
\item \textbf{[debian-wiki-swap]} Debian Wiki: Swap --- \textit{Debian Community} (Current) \textcolor{cOk}{[consult\'e]}\\ \url{https://wiki.debian.org/Swap}
\item \textbf{[debian-wiki-ext4]} Debian Wiki: Ext4 --- \textit{Debian Community} (Current) \textcolor{cOk}{[consult\'e]}\\ \url{https://wiki.debian.org/Ext4}
\item \textbf{[debian-wiki-fstab]} Debian Wiki: fstab --- \textit{Debian Community} (Current) \textcolor{cOk}{[consult\'e]}\\ \url{https://wiki.debian.org/fstab}
\item \textbf{[ubuntu-system-requirements]} Ubuntu System Requirements - Official Documentation --- \textit{Ubuntu/Canonical} (Current) \textcolor{cOk}{[consult\'e]}\\ \url{https://ubuntu.com/server/docs/reference/installation/system-requirements/}
\item \textbf{[microsoft-uefi-spec]} UEFI/GPT-based hard drive partitions --- \textit{Microsoft Learn} (Current) \textcolor{cOk}{[consult\'e]}\\ \url{https://learn.microsoft.com/en-us/windows-hardware/manufacture/desktop/configure-uefigpt-based-hard-drive-partitions}
\item \textbf{[ubuntu-cis-compliance]} CIS Compliance with Ubuntu - Official Ubuntu Security Documentation --- \textit{Canonical/Ubuntu} (Current) \textcolor{cOk}{[consult\'e]}\\ \url{https://ubuntu.com/security/certifications/docs/usg/cis}
\item \textbf{[ubuntu-zram-guide]} How to Use ZRAM on Ubuntu 24.04 --- \textit{UbuntuHandbook/Community} (2024) \textcolor{cOk}{[consult\'e]}\\ \url{https://ubuntuhandbook.org/index.php/2024/08/enable-zram-ubuntu/}
\end{itemize}
\subsection{SUSE/SLES}
\begin{itemize}[leftmargin=1.6em]
\item \textbf{[suse-sles12-sp5-deployment]} Advanced Disk Setup | Deployment Guide | SLES 12 SP5 --- \textit{SUSE LLC} (SLES 12 SP5) \textcolor{cOk}{[consult\'e]}\\ \url{https://documentation.suse.com/sles/12-SP5/html/SLES-all/cha-advdisk.html}
\item \textbf{[suse-sles12-sp5-expert-partitioner]} Expert Partitioner | Deployment Guide | SLES 12 SP5 --- \textit{SUSE LLC} (SLES 12 SP5) \textcolor{cOk}{[consult\'e]}\\ \url{https://documentation.suse.com/sles/12-SP5/html/SLES-all/cha-expert-partitioner.html}
\item \textbf{[suse-sles12-sp5-filesystems]} Overview of File Systems in Linux | Storage Administration Guide | SLES 12 SP5 --- \textit{SUSE LLC} (SLES 12 SP5) \textcolor{cOk}{[consult\'e]}\\ \url{https://documentation.suse.com/sles/12-SP5/html/SLES-all/cha-filesystems.html}
\item \textbf{[suse-sles12-sp5-uefi]} UEFI (Unified Extensible Firmware Interface) | Administration Guide | SLES 12 SP5 --- \textit{SUSE LLC} (SLES 12 SP5) \textcolor{cOk}{[consult\'e]}\\ \url{https://documentation.suse.com/sles/12-SP5/html/SLES-all/cha-uefi.html}
\item \textbf{[suse-sles15-sp7-expert-partitioner]} Expert Partitioner | Deployment Guide | SLES 15 SP7 --- \textit{SUSE LLC} (SLES 15 SP7) \textcolor{cOk}{[consult\'e]}\\ \url{https://documentation.suse.com/sles/15-SP7/html/SLES-all/cha-expert-partitioner.html}
\item \textbf{[suse-sles15-sp7-filesystems]} Overview of file systems in Linux | Storage Administration Guide | SLES 15 SP7 --- \textit{SUSE LLC} (SLES 15 SP7) \textcolor{cOk}{[consult\'e]}\\ \url{https://documentation.suse.com/sles/15-SP7/html/SLES-all/cha-filesystems.html}
\item \textbf{[suse-sles15-sp7-uefi]} UEFI (Unified Extensible Firmware Interface) | Administration Guide | SLES 15 SP5 --- \textit{SUSE LLC} (SLES 15 SP5) \textcolor{cOk}{[consult\'e]}\\ \url{https://documentation.suse.com/sles/15-SP5/html/SLES-all/cha-uefi.html}
\item \textbf{[cis-suse-benchmark]} CIS SUSE Linux Enterprise Server Benchmarks (v2.0.0-v2.1.0 for SLES 12; v1.0.0+ for SLES 15) --- \textit{Center for Internet Security (CIS)} (v2.1.0 (SLES 12), v1.0.0+ (SLES 15)) \textcolor{cMed}{[non v\'erifi\'e]}\\ \url{https://www.cisecurity.org/benchmark/suse_linux}
\item \textbf{[fhs-3.0]} Filesystem Hierarchy Standard (FHS) 3.0 --- \textit{Linux Foundation} (3.0 (June 2015)) \textcolor{cMed}{[non v\'erifi\'e]}\\ \url{https://refspecs.linuxfoundation.org/FHS_3.0/fhs-3.0.html}
\item \textbf{[openSUSE-SDB-BTRFS]} SDB:BTRFS --- \textit{openSUSE Project} (Current) \textcolor{cMed}{[non v\'erifi\'e]}\\ \url{https://en.opensuse.org/SDB:BTRFS}
\item \textbf{[openSUSE-SDB-Partitioning]} SDB:Partitioning --- \textit{openSUSE Project} (Current) \textcolor{cOk}{[consult\'e]}\\ \url{https://en.opensuse.org/SDB:Partitioning}
\end{itemize}
\subsection{AIX}
\begin{itemize}[leftmargin=1.6em]
\item \textbf{[ibm-release-notes-71]} AIX 7.1 Release Notes and Installation Tips --- \textit{IBM} (AIX 7.1 (7100-05 through 7.1 TL5)) \textcolor{cOk}{[consult\'e]}\\ \url{https://www.ibm.com/docs/en/aix/7.1.0?topic=notes-aix-715-release}
\item \textbf{[ibm-release-notes-72]} AIX Version 7.2 Release Notes --- \textit{IBM} (AIX 7.2.0 and Technology Levels) \textcolor{cOk}{[consult\'e]}\\ \url{https://www.ibm.com/docs/en/aix/7.2.0?topic=notes-aix-720-release}
\item \textbf{[ibm-release-notes-73]} AIX 7.3 Release Notes --- \textit{IBM} (AIX 7.3.0 and Technology Levels) \textcolor{cOk}{[consult\'e]}\\ \url{https://www.ibm.com/docs/en/aix/7.3.0?topic=notes-aix-73-release}
\item \textbf{[ibm-logical-volumes]} Logical volume storage and Volume groups concepts --- \textit{IBM} (AIX 7.1-7.3 Documentation) \textcolor{cMed}{[non v\'erifi\'e]}\\ \url{https://www.ibm.com/docs/en/aix/7.2.0?topic=management-logical-volume-storage}
\item \textbf{[ibm-paging-spaces]} Paging spaces placement and sizes --- \textit{IBM} (AIX 7.1-7.3 Performance Guidelines) \textcolor{cMed}{[non v\'erifi\'e]}\\ \url{https://www.ibm.com/docs/en/aix/7.2.0?topic=guidelines-paging-spaces-placement-sizes}
\item \textbf{[ibm-jfs2-log-device]} Increasing the Size of a JFS/JFS2 Log Device --- \textit{IBM Support} (AIX 7.1-7.3) \textcolor{cOk}{[consult\'e]}\\ \url{https://www.ibm.com/support/pages/increasing-size-jfsjfs2-log-device}
\item \textbf{[ibm-dump-devices]} Managing System Dump Devices --- \textit{IBM Support} (AIX 7.1-7.3) \textcolor{cOk}{[consult\'e]}\\ \url{https://www.ibm.com/support/pages/managing-system-dump-devices}
\item \textbf{[ibm-livedumps]} AIX Livedumps --- \textit{IBM Support} (AIX 6.1-7.3) \textcolor{cOk}{[consult\'e]}\\ \url{https://www.ibm.com/support/pages/aix-livedumps}
\item \textbf{[ibm-datavg-best-practice]} Give rootvg the space it needs and datavg best practices --- \textit{IBM Developer / IBM Support} (AIX 5L-7.3) \textcolor{cMed}{[non v\'erifi\'e]}\\ \url{https://www.ibm.com/developerworks/aix/library/au-rootvg/index.html}
\item \textbf{[ibm-redbook-lvm]} AIX Logical Volume Manager, From A to Z --- \textit{IBM Redbooks} (AIX 5L-7.x) \textcolor{cMed}{[non v\'erifi\'e]}\\ \url{https://www.redbooks.ibm.com/redbooks/pdfs/sg245432.pdf}
\item \textbf{[cis-aix-benchmark]} CIS AIX Benchmark v1.0.1 --- \textit{Center for Internet Security} (AIX 4.3.2-5.1 (legacy benchmark)) \textcolor{cMed}{[non v\'erifi\'e]}\\ \url{https://www.cisecurity.org/-/jssmedia/Project/cisecurity/cisecurity/data/media/files/uploads/2017/04/CIS_AIX_Benchmark_v101.pdf}
\item \textbf{[ibm-mount-command-72]} mount Command Documentation --- \textit{IBM} (AIX 7.2) \textcolor{cMed}{[non v\'erifi\'e]}\\ \url{https://www.ibm.com/support/knowledgecenter/ssw_aix_72/m_commands/mount.html}
\item \textbf{[ibm-mount-command-71]} mount Command Documentation --- \textit{IBM} (AIX 7.1) \textcolor{cMed}{[non v\'erifi\'e]}\\ \url{https://www.ibm.com/support/knowledgecenter/ssw_aix_71/com.ibm.aix.cmds3/mount.htm}
\end{itemize}
\subsection{Windows Server}
\begin{itemize}[leftmargin=1.6em]
\item \textbf{[ms-uefi-gpt-config]} UEFI/GPT-based hard drive partitions --- \textit{Microsoft Learn} (2022-07-01 (updated 2026-05-18)) \textcolor{cOk}{[consult\'e]}\\ \url{https://learn.microsoft.com/en-us/windows-hardware/manufacture/desktop/configure-uefigpt-based-hard-drive-partitions?view=windows-11}
\item \textbf{[ms-hardware-reqs]} Hardware Requirements for Windows Server --- \textit{Microsoft Learn} (2025-07-02 (updated 2026-02-16)) \textcolor{cOk}{[consult\'e]}\\ \url{https://learn.microsoft.com/en-us/windows-server/get-started/hardware-requirements}
\item \textbf{[ms-winre-partition]} Disk partition requirement for using Windows RE tools on a UEFI-based computer --- \textit{Microsoft Learn} (2026-02-12 (updated 2026-02-19)) \textcolor{cOk}{[consult\'e]}\\ \url{https://learn.microsoft.com/en-us/troubleshoot/windows-client/windows-security/disk-partition-requirement-use-windows-re-tool}
\item \textbf{[ms-bitlocker-planning]} BitLocker planning guide --- \textit{Microsoft Learn} (2025-07-29) \textcolor{cOk}{[consult\'e]}\\ \url{https://learn.microsoft.com/en-us/windows/security/operating-system-security/data-protection/bitlocker/planning-guide}
\item \textbf{[ms-refs-overview]} Resilient File System (ReFS) overview --- \textit{Microsoft Learn} (2025-07-22 (updated 2025-12-09)) \textcolor{cOk}{[consult\'e]}\\ \url{https://learn.microsoft.com/en-us/windows-server/storage/refs/refs-overview}
\item \textbf{[ms-volume-planning]} Plan volumes on Azure Local and Windows Server clusters --- \textit{Microsoft Learn} (2026-03-09 (updated 2026-04-15)) \textcolor{cOk}{[consult\'e]}\\ \url{https://learn.microsoft.com/en-us/windows-server/storage/storage-spaces/plan-volumes}
\item \textbf{[ms-pagefile-sizing]} How to determine the appropriate page file size for 64-bit versions of Windows --- \textit{Microsoft Learn} (2026-02-12 (updated 2026-02-19)) \textcolor{cOk}{[consult\'e]}\\ \url{https://learn.microsoft.com/en-us/troubleshoot/windows-client/performance/how-to-determine-the-appropriate-page-file-size-for-64-bit-versions-of-windows}
\item \textbf{[cis-benchmarks-portal]} CIS Microsoft Windows Server Benchmarks --- \textit{Center for Internet Security} (2025 (v4.0.0 for 2022, v1.0.0 for 2025)) \textcolor{cOk}{[consult\'e]}\\ \url{https://www.cisecurity.org/benchmark/microsoft_windows_server}
\item \textbf{[ms-storage-spaces-direct]} Storage Spaces Direct Hardware Requirements in Windows Server --- \textit{Microsoft Learn} (Current) \textcolor{cOk}{[consult\'e]}\\ \url{https://learn.microsoft.com/en-us/windows-server/storage/storage-spaces/storage-spaces-direct-hardware-requirements}
\end{itemize}
\subsection{Cross-distro baseline (FHS/CIS)}
\begin{itemize}[leftmargin=1.6em]
\item \textbf{[FHS3.0-ch03]} Filesystem Hierarchy Standard 3.0 - Chapter 3: The Root Filesystem --- \textit{Linux Foundation} (FHS 3.0 (published June 2015, republished by FreeDesktop November 2025)) \textcolor{cOk}{[consult\'e]}\\ \url{https://refspecs.linuxfoundation.org/FHS_3.0/fhs/ch03.html}
\item \textbf{[FHS3.0-ch04]} Filesystem Hierarchy Standard 3.0 - Chapter 4: The /usr Hierarchy --- \textit{Linux Foundation} (FHS 3.0) \textcolor{cOk}{[consult\'e]}\\ \url{https://refspecs.linuxfoundation.org/FHS_3.0/fhs-3.0.txt}
\item \textbf{[FHS3.0-ch05]} Filesystem Hierarchy Standard 3.0 - Chapter 5: The /var Hierarchy --- \textit{Linux Foundation} (FHS 3.0) \textcolor{cMed}{[non v\'erifi\'e]}\\ \url{https://refspecs.linuxfoundation.org/FHS_3.0/fhs/ch05.html}
\item \textbf{[FHS3.0-ch06]} Filesystem Hierarchy Standard 3.0 - /opt and /srv --- \textit{Linux Foundation} (FHS 3.0) \textcolor{cOk}{[consult\'e]}\\ \url{https://refspecs.linuxfoundation.org/FHS_3.0/fhs-3.0.html}
\item \textbf{[FHS3.0-ch07]} Filesystem Hierarchy Standard 3.0 - /srv Data Directory --- \textit{Linux Foundation} (FHS 3.0) \textcolor{cOk}{[consult\'e]}\\ \url{https://refspecs.linuxfoundation.org/FHS_3.0/fhs-3.0.html}
\item \textbf{[CIS-Linux-1.1.2]} CIS Benchmarks for Linux - Control 1.1.2: Ensure /tmp is a separate partition --- \textit{Center for Internet Security (CIS)} (CIS Ubuntu 20.04/22.04/24.04 LTS Benchmarks; CIS RHEL 8/9 Benchmarks) \textcolor{cOk}{[consult\'e]}\\ \url{https://www.cisecurity.org/benchmark/ubuntu_linux}
\item \textbf{[CIS-Linux-1.1.4]} CIS Benchmarks for Linux - Control 1.1.4: Ensure /var is a separate partition --- \textit{Center for Internet Security (CIS)} (CIS Linux Benchmarks (multi-distribution)) \textcolor{cOk}{[consult\'e]}\\ \url{https://www.cisecurity.org/benchmark/ubuntu_linux}
\item \textbf{[CIS-Linux-1.1.5]} CIS Benchmarks for Linux - Control 1.1.5: Ensure /var/log is a separate partition --- \textit{Center for Internet Security (CIS)} (CIS Linux Benchmarks) \textcolor{cOk}{[consult\'e]}\\ \url{https://www.cisecurity.org/benchmark/ubuntu_linux}
\item \textbf{[CIS-Linux-1.1.6]} CIS Benchmarks for Linux - Control 1.1.6: Ensure /var/log/audit is a separate partition --- \textit{Center for Internet Security (CIS)} (CIS Linux Benchmarks) \textcolor{cOk}{[consult\'e]}\\ \url{https://www.cisecurity.org/benchmark/ubuntu_linux}
\item \textbf{[CIS-Linux-1.1.7]} CIS Benchmarks for Linux - Control 1.1.7: Ensure /var/tmp is configured --- \textit{Center for Internet Security (CIS)} (CIS Linux Benchmarks) \textcolor{cOk}{[consult\'e]}\\ \url{https://www.cisecurity.org/benchmark/ubuntu_linux}
\item \textbf{[CIS-Linux-1.1.9]} CIS Benchmarks for Linux - Control 1.1.9: Ensure /home is a separate partition --- \textit{Center for Internet Security (CIS)} (CIS Linux Benchmarks) \textcolor{cOk}{[consult\'e]}\\ \url{https://www.cisecurity.org/benchmark/ubuntu_linux}
\item \textbf{[CIS-Linux-1.1.15]} CIS Benchmarks for Linux - Control 1.1.15-1.1.17: Ensure /dev/shm mount options --- \textit{Center for Internet Security (CIS)} (CIS Linux Benchmarks) \textcolor{cOk}{[consult\'e]}\\ \url{https://www.cisecurity.org/benchmark/ubuntu_linux}
\item \textbf{[CIS-Linux-partition]} CIS Benchmark - Linux Filesystem Partitions --- \textit{Center for Internet Security} (2024-2026) \textcolor{cMed}{[non v\'erifi\'e]}\\ \url{https://breadandwater.dev/cis-benchmark-linux-filesystem-partitions/}
\item \textbf{[CIS-Linux-hardening]} CIS Benchmarks for Linux - Filesystem Configuration Controls --- \textit{Center for Internet Security} (Latest (Ubuntu/RHEL/Debian)) \textcolor{cOk}{[consult\'e]}\\ \url{https://www.cisecurity.org/benchmark/ubuntu_linux}
\item \textbf{[RHEL-install-guide]} Red Hat Enterprise Linux 8, 9, 10 - Installation Guide: Partitioning Reference --- \textit{Red Hat} (RHEL 8/9/10 (2024-2026)) \textcolor{cMed}{[non v\'erifi\'e]}\\ \url{https://docs.redhat.com/en/documentation/red_hat_enterprise_linux/10/html/managing_storage_devices/disk-partitions}
\item \textbf{[RHEL-filesystem-guide]} Red Hat Enterprise Linux 8, 9 - Managing File Systems Guide --- \textit{Red Hat} (RHEL 8/9 (2022-2026)) \textcolor{cMed}{[non v\'erifi\'e]}\\ \url{https://docs.redhat.com/en/documentation/red_hat_enterprise_linux/9/html/managing_file_systems/overview-of-available-file-systems_managing-file-systems}
\item \textbf{[RHEL-swap-guide]} Red Hat Enterprise Linux - Getting Started with Swap --- \textit{Red Hat} (RHEL 8/9/10) \textcolor{cOk}{[consult\'e]}\\ \url{https://docs.redhat.com/en/documentation/red_hat_enterprise_linux/10/html/managing_storage_devices/getting-started-with-swap}
\item \textbf{[Debian-install-guide]} Debian Linux Installation Guide - Partitioning Schemes --- \textit{Debian Project} (Debian Stable (Bookworm/Trixie)) \textcolor{cOk}{[consult\'e]}\\ \url{https://www.debian.org/releases/stable/arm64/apcs03.en.html}
\item \textbf{[Debian-secure-manual]} Securing Debian Manual - Partitioning the System --- \textit{Debian Project} (Debian Security Manual) \textcolor{cOk}{[consult\'e]}\\ \url{https://www.debian.org/doc/manuals/securing-debian-manual/ch03s02.en.html}
\item \textbf{[Ubuntu-partitioning]} Ubuntu Community Help Wiki - Partitioning Schemes --- \textit{Canonical/Ubuntu Community} (Ubuntu 20.04 LTS - 24.04 LTS) \textcolor{cOk}{[consult\'e]}\\ \url{https://help.ubuntu.com/community/PartitioningSchemes}
\item \textbf{[CIS-Ubuntu]} CIS Ubuntu Linux Benchmarks --- \textit{Center for Internet Security} (Ubuntu 20.04/22.04/24.04 LTS) \textcolor{cOk}{[consult\'e]}\\ \url{https://ubuntu.com/security/certifications/docs/16-18/cis/customization}
\item \textbf{[SUSE-SLES-expert-partitioner]} SUSE Linux Enterprise Server 15 SP7 - Expert Partitioner --- \textit{SUSE} (SLES 15 SP6/SP7 (2022-2026)) \textcolor{cOk}{[consult\'e]}\\ \url{https://documentation.suse.com/sles/15-SP7/html/SLES-all/cha-expert-partitioner.html}
\item \textbf{[TechGirlKB]} How to Create CIS-Compliant Partitions on AWS --- \textit{TechGirlKB (practitioner guide)} (2019-2026) \textcolor{cOk}{[consult\'e]}\\ \url{https://techgirlkb.guru/2019/08/how-to-create-cis-compliant-partitions-on-aws/}
\item \textbf{[Proxmox-CIS]} CIS Benchmark Partition Configuration Forum --- \textit{Proxmox Community} (2021-2026) \textcolor{cOk}{[consult\'e]}\\ \url{https://forum.proxmox.com/threads/cis-benchmark-fail-tmp-var-var-tmp-var-log-var-log-audit-and-home-parition-configuration.95033/}
\end{itemize}
\section{Limites \& points à v\'erifier}
\noindent\textbf{Qualit\'e des sources :} ["Pervasive retrieved=false on the load-bearing sources. In the RHEL block, 13 of 17 sources are retrieved=false - including ALL three CIS RHEL benchmarks and the RHEL 8/9 install guides that supply the actual size/mount figures. coverage\_notes admits content was 'behind authentication walls' and 'confirmed via web search', i.e., values are second-hand, not directly verified. This is the single biggest quality issue.", "All CIS benchmark sources across every block are retrieved=false, yet CIS is the cited authority for nearly all nodev/nosuid/noexec mount options and 'mandatory separate partition' claims. The hardening recommendations therefore rest on unretrieved primary sources.", "Non-authoritative / dubious-host URLs used for primary standards: cis-debian-11 and cis-debian-12 point to itsecure.hu (third-party WordPress upload), cis-ubuntu-20 to edu.anarcho-copy.org, and cis-ubuntu-24 to scribd.com. These are copyright-infringing re-hosts, not CIS - replace with cisecurity.org/workbench citations. Same problem with CIS-Linux-partition (breadandwater.dev), TechGirlKB, and Proxmox-CIS forum used as CIS evidence in the baseline block.", "Source-ID references that don't exist in the block's sources array (broken citations): RHEL7 partitions cite 'cis-rhel7' / RHEL8 cite 'cis-rhel8' / RHEL9 cite 'cis-rhel9' and 'rhel7-storage-admin' etc. - those IDs ARE defined, but the cross-distro baseline partitions cite 'Debian-secure-manual', 'Ubuntu-partitioning', 'CIS-Ubuntu', 'Debian-install-guide', 'RHEL-install-guide', 'RHEL-filesystem-guide' which exist, however baseline partition 'TechGirlKB' is cited and defined but is a personal blog (not authoritative for audit-partition sizing). Verify every source\_id resolves before PDF build.", "Stale-dating / future-dating anomalies that undermine trust: many Windows Server sources carry 'updated 2026-05-18 / 2026-04-15 / 2026-02-19' dates that are AFTER today's date (2026-06-12 is fine, but 2026-05/04 'updated' on Microsoft Learn should be sanity-checked) and FHS is described as 'republished by FreeDesktop November 2025' - verify this claim, it is unusual and uncited to a retrievable page.", "Mixed/!inconsistent CIS versioning: cis-rhel8 listed as 'v2.0.0' (no date) while cis-rhel9 is 'v2.0.0 (2024-06-24)' and cis-rhel7 is 'v4.0.0' - RHEL7 at v4.0.0 while RHEL8/9 at v2.0.0 is suspicious (benchmark major versions are per-distro but the jump should be verified). cis-suse-benchmark bundles multiple versions into one ID, defeating traceability.", "ibm-mount-command and several IBM logical-volume/paging sources are retrieved=false with note 'some IBM docs pages return 403'; the AIX mount-option and paging-formula claims thus rest on unverified pages. ibm-datavg-best-practice URL is a developerWorks link (developerWorks was retired) - likely dead, should be re-pointed."]\par\medskip
\noindent\textbf{Lacunes / manques :}
\begin{itemize}[leftmargin=1.6em]
\item RHEL family: no RHEL 6 (still in ELS) and no RHEL 10 (GA'd 2025) even though the cross-distro block references 'RHEL 8, 9, 10'. Inconsistent scoping between blocks.
\item RHEL per-version blocks omit /dev/shm, /run, /srv, /usr, /opt, and /usr/local entirely, while the cross-distro baseline lists them as CIS/FHS-relevant. /dev/shm in particular is a CIS-mandated mount-options control and should appear in every RHEL version.
\item RHEL: RHEL-as-base rebuilds (Rocky/Alma/Oracle) are only mentioned in coverage\_notes prose, not modeled as versions; if the PDF claims to cover them, add explicit entries or a clear note.
\item Debian/Ubuntu: missing /run, /srv, /opt, /usr/local mount points that appear in the cross-distro baseline. /var/log/audit is also absent from the Debian/Ubuntu per-version partition lists despite being a core CIS control (present in RHEL and baseline blocks).
\item Debian/Ubuntu: Ubuntu 18.04 LTS (still ESM) and the forthcoming Debian 13 (Trixie) are not covered; coverage\_notes says 'five most recent' but scope choice should be stated explicitly in the PDF.
\item Ubuntu modern default is a swapfile (and zram on 24.04), but every Ubuntu version still lists swap requirement as 'mandatory' separate\_filesystem 'yes' - the default-install reality (swapfile, no swap partition) is a gap/contradiction.
\item SUSE/SLES: SP-level granularity is collapsed ('SLES 15 (SP1-SP7)'). Boot-LV / EFI minimums and Btrfs subvolume layout changed across SP releases; lumping SP1 through SP7 hides version-specific differences. SLES 12 SP5 is EOL (general support ended) - note its support status.
\item SUSE/SLES: no per-version `coverage\_notes` URL set verified - the cis-suse-benchmark and openSUSE sources are all retrieved=false; /dev/shm and /var/tmp recommendations are explicitly described as 'inferred from general CIS Linux principles', a real evidentiary gap.
\item AIX: no AIX 7.1 vs 7.2 vs 7.3 difference beyond hd5 boot size; CIS AIX benchmark cited is legacy (4.3.2-5.1) and retrieved=false, so mount-option hardening (rw only) is effectively unsupported for 7.x. Note that JFS2 mount options (e.g., nodev/nosuid) are not modeled at all.
\item Windows Server: no separate /pagefile or recovery-image partition sizing detail beyond prose; CIS Windows benchmarks cited but no partition-specific controls extracted. Windows Server 2012 R2 (still ESU) not covered - acceptable if scope is LTSC 2016+, but state it.
\item Cross-distro baseline: contains a third RHEL representation ('RHEL 8, 9, 10 - Vendor Guidance') with DIFFERENT figures than the dedicated RHEL block (e.g., / min 15 GiB vs 5 GiB; /boot max 1 GiB; /var/log 5 GiB). These two models should be reconciled or the PDF should explain why they differ.
\end{itemize}
\noindent\textbf{Chiffres à re-v\'erifier :}
\begin{itemize}[leftmargin=1.6em]
\item RHEL 7 swap max\_size '128 GiB' as a hard partition max is misleading: the 10\%/128 GB cap applies only to AUTOMATIC partitioning swap sizing, not a filesystem maximum. Flag 'verify' - presenting it as the partition max\_size mis-states the source.
\item RHEL 7 / min\_size '3.0 GiB': RHEL 7 install guide root minimum is commonly cited as 10 GiB recommended / smaller technically allowed; the 3.0 GiB figure (and the 'XFS' fstype on a row whose notes say /boot must be ext4) is internally muddled. The RHEL7 / row says fstype XFS but default\_fstype text and notes claim 'Default filesystem for / and /boot is ext4' - contradictory; verify.
\item RHEL 7 /boot min\_size '250 MiB' coexists with notes saying default 1 GB since 7.3 and the swap/boot narrative - the 250 MiB pre-7.3 minimum vs 1 GiB default should be reconciled; flag the 250 MiB row.
\item Cross-distro 'RHEL 8/9/10 Vendor Guidance' / min\_size '15 GiB' directly contradicts the dedicated RHEL 8/9 blocks' '5.0 GiB'. Both cannot be the authoritative minimum - verify which is min vs recommended.
\item Ubuntu 20.04 /boot/efi note 'creates minimum 538MiB' is an oddly specific, non-standard ESP size (Ubuntu's installer ESP is typically 512 MiB / 100-300 MiB on some images). 538 MiB looks wrong - flag 'verify'.
\item Debian/Ubuntu swap max\_size '2x RAM' combined with swap\_guidance 'maximum 2x RAM' contradicts the same text's modern 'round(sqrt(RAM))' guidance and Ubuntu's actual default 1-2 GB swapfile; the 2x RAM max is an outdated rule presented as current - flag.
\item AIX dump 'min\_size 1.5x RAM' as a partition minimum is a heuristic, not a documented hard minimum; and lg\_dumplv only exists for >=4 GB systems, so labeling dump 'mandatory / 1.5x RAM' universally is suspect - verify.
\item AIX hd6 paging 'max\_size 64 GB (per LV)' and the '512 MB + (RAM-256MB)*1.25' formula are old AIX 5L-era guidance; confirm still stated verbatim in 7.1-7.3 docs (sources for the formula are retrieved=false). Flag.
\item Windows Server WinRE free-space progression (52 MB -> 100 MB -> 200 MB across 2016/2022/2025) mixes 'WinRE partition min 300 MB' with 'free space' numbers; the 52 MB / 100 MB / 200 MB figures are easily conflated and at least one (Server 2016 '52 MB') is widely reported as 'no fixed minimum' - flag 'verify'.
\item Windows pagefile 'System-managed ... minimum to 3x RAM or 4 GB maximum (whichever larger)' garbles Microsoft's actual rule (min RAM/8 capped 32 GB; max 3x RAM or 4 GB whichever LARGER) - the phrasing conflates min and max; verify wording.
\item SUSE /boot/efi 'min 260 MiB' is unusual (ESP minimums are normally 100-260 MiB depending on cluster size); 260 MiB is defensible for FAT32 but the max '500 MiB' as a standard cap is arbitrary - low priority but flag.
\item Cross-distro baseline /boot/efi min '50 MiB' max '150 MiB' contradicts every per-distro block (which use 200-260 MiB minimum). 50 MiB ESP is too small for modern UEFI and conflicts with the rest of the compilation - flag prominently.
\end{itemize}
\par\medskip\noindent\small\textit{"The compilation is broad (6 OS families, \textasciitilde{}20 versions) and well-structured, but it is NOT yet PDF-ready because (1) the overwhelming majority of load-bearing sources are retrieved=false - including ALL CIS benchmarks and most RHEL/AIX primary docs - so most figures are second-hand and unverifiable as written; (2) several CIS sources point to copyright-infringing third-party re-hosts (itsecure.hu, anarcho-copy, scribd) instead of cisecurity.org, which is both a quality and a legitimacy problem; and (3) there are direct internal contradictions, most notably two different RHEL models (dedicated block / min 5 GiB vs cross-distro 'vendor guidance' / min 15 GiB) and an EFI minimum that ranges from 50 MiB to 260 MiB across blocks. Priority actions before publishing: reconcile the duplicate RHEL models and the /boot/efi minimums; re-source every CIS citation to the official CIS workbench and actually retrieve at least the size/mount-option controls; mark the swap 'max 2x RAM', RHEL swap '128 GiB max', Windows WinRE free-space numbers, AIX dump/paging formulas, and Ubuntu '538 MiB ESP' rows as 'verify'; and add the missing /dev/shm, /run, /srv, /var/log/audit rows to the Debian/Ubuntu and RHEL per-version tables (or state explicitly that they are out of scope). Coverage gaps (RHEL 10, Debian 13, SLES SP granularity, RHEL-rebuilds) should be either filled or scoped in writing."}
\end{document}